\documentclass[11pt]{amsart}
\usepackage{amsmath,amssymb,amsthm}
\usepackage{url}
\usepackage{tikz}
\usetikzlibrary{calc}
\setlength{\emergencystretch}{3em}

\newtheorem{theorem}{Theorem}[section]
\newtheorem{proposition}[theorem]{Proposition}
\newtheorem{lemma}[theorem]{Lemma}
\newtheorem{corollary}[theorem]{Corollary}
\newtheorem*{theoremA}{Theorem A}
\newtheorem*{theoremB}{Theorem B}
\newtheorem*{theoremC}{Theorem C}
\theoremstyle{definition}
\newtheorem{definition}[theorem]{Definition}
\newtheorem{example}[theorem]{Example}
\theoremstyle{remark}
\newtheorem{remark}[theorem]{Remark}

\newcommand{\Q}{\mathbb Q}
\newcommand{\Z}{\mathbb Z}
\newcommand{\C}{\mathbb C}
\newcommand{\R}{\mathbb R}
\newcommand{\PP}{\mathbb P}
\newcommand{\End}{\operatorname{End}}
\newcommand{\Hom}{\operatorname{Hom}}
\newcommand{\Nm}{\operatorname{Nm}}
\newcommand{\rk}{\operatorname{rk}}
\newcommand{\Tr}{\operatorname{Tr}}
\newcommand{\Ind}{\operatorname{Ind}}
\newcommand{\Hg}{\operatorname{Hg}}
\newcommand{\disc}{\operatorname{disc}}
\newcommand{\age}{\operatorname{age}}
\newcommand{\cH}{\mathrm{Hur}}
\newcommand{\cHn}{\mathcal H_n}
\newcommand{\Vp}{V_{\rm par}}
\newcommand{\VpC}{V_{{\rm par},\C}}
\newcommand{\cM}{\mathcal M}
\newcommand{\cD}{\mathcal D}
\newcommand{\fd}{\mathfrak d}

\title[Weil classes for $\Q(\sqrt{-3})$]{Weil classes on abelian sixfolds and eightfolds of Weil type for $\Q(\sqrt{-3})$}
\author{Benjamin Severn}
\email{benzsevern@gmail.com}
\subjclass[2020]{Primary 14C30; Secondary 14K20, 14D07, 14H30}
\keywords{Hodge conjecture, Weil classes, abelian varieties of Weil type, Prym--Tyurin varieties, complex reflection groups}
\date{21 October 2026 (version 1.5)}

\begin{document}

\begin{abstract}
Let $K=\Q(\sqrt{-3})$. We prove that the Weil classes of every polarized complex abelian eightfold of Weil type for $K$ with split Hermitian discriminant are algebraic, and deduce the same for every abelian sixfold of Weil type for $K$, of every discriminant. For sixfolds whose Hodge group is the full special unitary group, the Hodge conjecture follows on all powers. The eightfolds are Prym--Tyurin varieties of Galois covers of the projective line with monodromy group the Shephard--Todd group $G_{34}$, branched at twenty points with reflections as local monodromy. Algebraic cycles come from Schoen's construction on cyclic subcovers and reach the Prym--Tyurin factor through Milne's theorem on Lefschetz classes; a character identity that makes this possible follows from Brauer's induction theorem. The family dominates the sixteen-dimensional moduli space of split eightfolds; we prove this from exact data at a maximally degenerate point of the Hurwitz space, by the nilpotent orbit theorem. The sixfolds are reached by multiplying with a product of CM elliptic curves of suitable polarization, as in Schoen's product argument.
\end{abstract}
\maketitle

%% ------------------------------------------------------------------
\section{Introduction}
%% ------------------------------------------------------------------

\subsection{Weil classes}
Let $A$ be a complex abelian variety of dimension $2n$ with an action $\iota\colon K\to\End(A)\otimes\Q$ of an imaginary quadratic field $K$, such that the two embeddings of $K$ in $\C$ each occur with multiplicity $n$ on the tangent space. Then the $K$-line
\[
W_K(A)=\textstyle\bigwedge^{2n}_KH^1(A,\Q)\subset H^{2n}(A,\Q)
\]
is a two-dimensional space of Hodge classes of type $(n,n)$, the \emph{Weil classes} of $A$ \cite{Weil}. For a general such $A$ they are not generated by divisor classes, so they are a natural test case for the Hodge conjecture \cite[\S5]{vanGeemen}. A polarization compatible with $\iota$ determines a $K$-Hermitian form $H$ on $H_1(A,\Q)$ of signature $(n,n)$, and its discriminant $\disc H\in\Q^\times/\Nm K^\times$ is an invariant of the polarized rational Hodge structure. The polarized abelian $2n$-folds of Weil type for $K$ with a fixed discriminant form one family, parametrized by a quotient of the Hermitian symmetric domain of $SU(n,n)$, of dimension $n^2$; we call its members the abelian $2n$-folds of Weil type for $K$ with that discriminant. The discriminant $(-1)^n$ is called \emph{split}: it is the discriminant of the Hermitian forms that have an $n$-dimensional totally isotropic subspace.

\subsection{What was known}
For fourfolds, the Weil classes are algebraic for every imaginary quadratic field and every discriminant: Markman proved the split case \cite{Markman19}, and deduced the general case from split sixfolds \cite{Markman25}. For sixfolds, the split case is known: Schoen for $\Q(\sqrt{-3})$ \cite{Schoen98}, Koike for $\Q(i)$ \cite{Koike}, and Markman for every $K$ \cite[Theorem~1.5.1]{Markman25}. Beyond the split class, the only sixfold case we know of is the class $-2$ for $\Q(\sqrt{-3})$ \cite{SevernG32}; Mostaed describes the non-split sixfold cases as open \cite[Introduction]{Mostaed}. Voisin's Bourbaki expos\'e \cite{Voisin} surveys this work.

In dimension eight and above, no full moduli space of abelian varieties of Weil type was previously known to have algebraic Weil classes. Markman's secant-sheaf method is carried out in dimension at most six, and he leaves dimension eight dependent on an open question \cite[\S\S1, 12]{Markman25b}. Van Geemen and Verra \cite{vGV} prove that the Weil classes are algebraic on the general member of a six-dimensional family of eightfolds of quaternion type, which lies inside a Weil family for $\Q(i)$ but does not fill it.

If split eightfolds were known for every $K$, the product argument below would give every Weil sixfold for every $K$. That is not known, and this is why the construction with $G_{34}$ is needed.

\subsection{Results}

\begin{theoremA}
Let $K=\Q(\sqrt{-3})$. The Weil classes of every polarized abelian eightfold of Weil type for $K$ with split discriminant are algebraic.
\end{theoremA}

\begin{theoremB}
Let $K=\Q(\sqrt{-3})$. The Weil classes of every polarized abelian sixfold of Weil type for $K$, of any discriminant, are algebraic.
\end{theoremB}

\begin{theoremC}
Let $A$ be an abelian sixfold as in Theorem~B whose Hodge group is $SU_H$. Then every Hodge class on every power $A^m$ is algebraic.
\end{theoremC}

Theorem~C holds for every sixfold of Weil type for $K$ outside a countable union of proper closed subvarieties of its moduli space. Section~\ref{sec:typeIII} gives a consequence for abelian sixfolds of type III.

\subsection{The construction in one page}\label{subsec:map}
The sixfolds are reached through eightfolds and a product construction:
\[
\begin{array}{c}
\text{$G_{34}$-covers $C\to\PP^1$}\\
\text{branched at $20$ points}
\end{array}
\ \Longrightarrow\
\begin{array}{c}
\text{Prym--Tyurin factor}\\
A_\sigma(C)\subset J(C),\ \dim 8
\end{array}
\ \Longrightarrow
\]
\[
\Longrightarrow\
\begin{array}{c}
\text{every split}\\
\text{eightfold}
\end{array}
\ \Longrightarrow\
\begin{array}{c}
A\times E\times\bar E\\
\text{for every sixfold }A
\end{array}
\]

\emph{The eightfolds.} The Shephard--Todd group $G_{34}$ is a complex reflection group of order $39\,191\,040$, acting on $K^6$ through a representation $\sigma$ whose character field is $K$. Take Galois covers $C\to\PP^1$ with group $G_{34}$, branched over twenty points, with a reflection as the local monodromy at each. The representation $\sigma$ occurs in $H^1(C)$ with multiplicity $8$, and the multiplicity space is the first cohomology of an abelian eightfold $A_\sigma(C)$, cut out of the Jacobian $J(C)$ by a rank-one idempotent of $\Q[G_{34}]$; the $\sigma$-isotypic part of $J(C)$ is isogenous to $A_\sigma(C)^6$. The centre of $\Q[G_{34}]$ acting on $\sigma$ is $K$, and it makes $A_\sigma(C)$ an eightfold of Weil type for $K$ (Section~\ref{sec:covers}). As $C$ moves in its $17$-dimensional Hurwitz space, $A_\sigma(C)$ moves in the $16$-dimensional moduli space of split eightfolds of Weil type for $K$.

\emph{Cycles.} For a subgroup $H'\le G_{34}$ and a character $\chi\colon H'\to\mu_m$, the cyclic cover $C/\ker\chi\to C/H'$ carries algebraic Weil classes on its Prym variety, by Schoen's theorem \cite{Schoen88}, under a condition on its branch data. These Prym varieties are factors of $J(C)$ cut out by idempotents of $\Q[G_{34}]$. A product of their Weil classes, contracted against a suitable class on a product of factors of $J(C)^R$, gives a nonzero Weil class on $A_\sigma(C)$. The contracting class is algebraic by Milne's theorem on Lefschetz classes \cite{Milne}, provided it is invariant under the centralizer of $G_{34}$; that is a condition on characters of a torus, an integral identity in its character group (Section~\ref{sec:supply}). For $G_{34}$ with reflections, the identity follows from Brauer's induction theorem, because every reflection is conjugate to its inverse (Section~\ref{sec:brauer}). So the Weil classes of $A_\sigma(C)$ are algebraic for every $C$.

\emph{Dominance and spreading.} If the family $A_\sigma(C)$ dominates the moduli space, the set of moduli points with algebraic Weil classes contains an open set; it is a countable union of closed subvarieties, so it is everything (Section~\ref{sec:spread}). Dominance is proved at a maximal degeneration of the twenty branch points. Let $\Vp\cong H_1(A_\sigma(C),\Q)$ be the eight-dimensional $K$-space of the family, with its Hermitian form; in Section~\ref{sec:herm} it is realized as a parabolic cohomology group. The logarithms of the local monodromies of the family there span the full $16$-dimensional abelian Lie algebra $\mathfrak n_I$ of endomorphisms of $\Vp$ that map $\Vp$ into a rational Lagrangian $I$ and kill $I$, and the nilpotent orbit theorem turns this into full rank of the period map nearby (Section~\ref{sec:boundary}). The required data, seventeen nilpotent matrices over $K$, are computed exactly from the monodromy tuple, which we print in Appendix~\ref{app:data}. The same data show that the Hermitian form of the family has a rational Lagrangian, so the family consists of split eightfolds.

\emph{From eightfolds to sixfolds.} If $A$ is a sixfold of Weil type for $K$ of discriminant $t$, and $E$ is an elliptic curve with complex multiplication by $\mathcal O_K$, then $A\times E\times\bar E$, with a product polarization of suitable degree on $\bar E$, is a split eightfold of Weil type for $K$. Its Weil classes are algebraic by Theorem~A, and contracting with the algebraic Weil classes of $E\times\bar E$ gives those of $A$ (Section~\ref{sec:compensation}). This is Schoen's product argument \cite[Proposition~10]{Schoen98}, by which he passes from sixfolds to fourfolds, in the dimension-free form of Voisin \cite[Lemme~2.9]{Voisin}; our only addition is the explicit choice of the compensating factor.

\emph{What is new.} The ingredients are classical. Schoen constructed the cycles on cyclic covers \cite{Schoen88}, and van Geemen and Verra already carried them to a non-abelian Prym factor through a cyclic intermediate quotient \cite[Corollary~4.10]{vGV}; dominance followed by specialization is the method of Schoen and Koike \cite{Schoen98,Koike}; and Alexeev, Donagi, Farkas, Izadi and Ortega proved dominance of a Prym--Tyurin map onto $\mathcal A_6$ at a maximally degenerate point of $\overline{\mathcal M}_{0,24}$ \cite{ADFIO}. Two things are new. The representation need not be induced from one cyclic subgroup: Weil classes are assembled from several cyclic covers by an integral identity in a character group, which Brauer's induction theorem supplies whenever the branch classes are real. And dominance of a Weil family is proved at a maximal degeneration, by exact linear algebra over $K$ on the monodromy tuple. The degeneration produces a rational Lagrangian, so this criterion only reaches the split discriminant (Proposition~\ref{prop:boundary}); that is why the method gives split eightfolds and then needs the product argument.

\emph{Inputs.} The proof uses Schoen's theorem on cyclic covers \cite{Schoen88}, Milne's theorem on Lefschetz classes \cite{Milne}, Brauer's induction theorem, the several-variable nilpotent orbit theorem \cite{Schmid,CKS}, the description of monodromy on parabolic cohomology by Dettweiler and Wewers \cite{DW1}, and a computation over $K$ whose output, a Gram matrix and seventeen $4\times4$ matrices, is printed in Appendix~\ref{app:data} together with the basis it refers to. Conditions (H1)--(H4) can be read off the printed matrices, and the matrices can be recomputed from the printed tuple and basis by exact linear algebra in dimension $8$ (Subsection~\ref{subsec:twists}).

\subsection{Organization}
Section~\ref{sec:weil} recalls abelian varieties of Weil type and their moduli. Section~\ref{sec:covers} constructs the Prym--Tyurin factor of a Galois cover. Section~\ref{sec:supply} produces algebraic Weil classes on it, and Section~\ref{sec:brauer} verifies the needed character identity for data with real local monodromy. Section~\ref{sec:spread} spreads algebraicity from a dominant family to the whole moduli space. Section~\ref{sec:herm} identifies the Hermitian form of the family with a parabolic cup product. Section~\ref{sec:boundary} proves the dominance criterion at a maximal degeneration. Section~\ref{sec:g34} applies everything to $G_{34}$ and proves Theorem~A, Section~\ref{sec:compensation} proves Theorem~B, and Section~\ref{sec:powers} proves Theorem~C. Section~\ref{sec:typeIII} treats abelian sixfolds of type III. Appendix~\ref{app:data} gives the data of the computation.

%% ------------------------------------------------------------------
\section{Abelian varieties of Weil type}\label{sec:weil}
%% ------------------------------------------------------------------

Throughout, $K$ is an imaginary quadratic field and $\vartheta\in K$ is a nonzero element with $\bar\vartheta=-\vartheta$.

\begin{definition}\label{def:weiltype}
A \emph{polarized abelian $2n$-fold of Weil type for $K$} is a triple $(A,\iota,E)$: a complex abelian variety $A$ of dimension $2n$, an embedding $\iota\colon K\to\End(A)\otimes\Q$ such that both embeddings of $K$ occur with multiplicity $n$ on $T_0A$, and a polarization $E$ (an alternating form on $H_1(A,\Q)$) with $E(\iota(a)x,y)=E(x,\iota(\bar a)y)$.
\end{definition}

Following van Geemen \cite[Lemma~5.2]{vanGeemen}, the form $H(x,y)=E(x,\iota(\vartheta)y)+\vartheta E(x,y)$ is $K$-Hermitian of signature $(n,n)$ on the $2n$-dimensional $K$-space $H_1(A,\Q)$. Its discriminant
\[
\disc H\in\Q^\times/\Nm K^\times
\]
does not depend on $\vartheta$, and it is unchanged by a $K$-linear isogeny that multiplies the polarization by a positive rational number, because such a change multiplies $\det H$ by a $2n$-th power of a rational number, which is a norm. The sign of $\disc H$ is $(-1)^n$. By Landherr's theorem \cite[\S5.4]{vanGeemen}, a nondegenerate $K$-Hermitian space is determined up to isometry by its dimension, signature and discriminant; in particular the Hermitian spaces of signature $(n,n)$ with a totally isotropic subspace of dimension $n$ are exactly those of discriminant $(-1)^n$, which we call \emph{split}.

\subsection{Moduli}\label{subsec:moduli}
Fix a $K$-Hermitian space $(V,H)$ of dimension $2n$, signature $(n,n)$ and discriminant $\fd$, and an $\mathcal O_K$-stable lattice $\Lambda\subset V$ on which the alternating form $\Tr_{K/\Q}(H/2\vartheta)$, scaled by a positive rational, is integral. The space $\cHn$ of $n$-dimensional complex subspaces of $V_\R$ (for the complex structure given by $\iota(\vartheta)/|\vartheta|$) on which $H$ is positive is the Hermitian symmetric domain of $SU(n,n)$, connected of dimension $n^2$. A point $V_+\in\cHn$ determines a complex structure $J'$ on $V_\R$ commuting with $K$, hence a Hodge structure of weight one on $V$ of Weil type for $K$, polarized by $\Tr_{K/\Q}(H/2\vartheta)$; and every such polarized Hodge structure arises in this way \cite[5.10 and the proof of 6.11]{vanGeemen}. Write $SU_H$ for the special unitary group of $(V,H)$, an algebraic group over $\Q$; it is the Hodge group of a very general member of the family below \cite[Theorem~6.11]{vanGeemen}. For $\ell\ge3$ the congruence subgroup $\Gamma(\ell)$ of the unitary group of $(\Lambda,H)$ is torsion free, and $\cM=\Gamma(\ell)\backslash\cHn$ is a smooth connected quasi-projective variety with a universal abelian scheme $f\colon\mathcal A\to\cM$ with $K$-action \cite[\S\S5.5--5.11]{vanGeemen}.

\begin{lemma}[Isogeny invariance]\label{lem:isogeny}
If $\psi\colon A'\to A$ is a $K$-linear isogeny, then $W_K(A)$ is spanned by algebraic classes if and only if $W_K(A')$ is. Every polarized abelian $2n$-fold of Weil type for $K$ with discriminant $\fd$ is $K$-isogenous to a fibre of $\mathcal A\to\cM$.
\end{lemma}

\begin{proof}
Pullback and pushforward by $\psi$ are algebraic correspondences, and $\psi^*$ maps $W_K(A)$ isomorphically onto $W_K(A')$, since $\psi^*$ is a $K$-linear isomorphism on $H^1(-,\Q)$. For the second statement, Landherr's theorem gives a $K$-isometry from $(H_1(A,\Q),H)$ to $(V,H)$; the image of $H_1(A,\Z)$ is commensurable with $\Lambda$, so a multiple of it lies in $\Lambda$, and multiplication by that integer is a $K$-linear isogeny from $A$ to the fibre of $\mathcal A$ at the point given by the complex structure of $A$.
\end{proof}

%% ------------------------------------------------------------------
\section{Prym--Tyurin factors of Galois covers}\label{sec:covers}
%% ------------------------------------------------------------------

\subsection{Hurwitz spaces}
Fix a finite group $G$ and conjugacy classes $\mathcal C_1,\dots,\mathcal C_r$ of $G$, $r\ge4$. The \emph{Nielsen class} is the set of tuples $(g_1,\dots,g_r)$ with $g_j\in\mathcal C_j$, $g_1\cdots g_r=1$, and $\langle g_1,\dots,g_r\rangle=G$. A connected Galois cover $C\to\PP^1$ with group $G$, branched over $r$ ordered points $p_1,\dots,p_r$, determines such a tuple up to simultaneous conjugation, once standard generators of $\pi_1(\PP^1-\{p_j\})$ are chosen: $g_j$ is the monodromy along a loop around $p_j$. The Hurwitz space $\cH$ of such covers is finite and \'etale over the configuration space of $r$ distinct ordered points; its connected components are smooth of dimension $r-3$ modulo $PGL_2$, and correspond to orbits of the pure braid group on the Nielsen class modulo conjugation \cite{FriedVolklein}.

For an irreducible complex representation $\tau$ of $G$ write $H^1(C,\C)=\bigoplus_\tau\tau\otimes M_\tau$ (Chevalley--Weil), and $m_\tau=\dim M_\tau$, the multiplicity of $\tau$ in $H^1(C,\C)$. The space $M_\tau$ carries a Hodge structure of weight one. For $g\in G$ write $c_\tau(g)=\rk(\tau(g)-1)$, and $\age\tau(g)=\sum_k\alpha_k$ when the eigenvalues of $\tau(g)$ are $e^{2\pi i\alpha_k}$ with $0\le\alpha_k<1$.

\begin{proposition}\label{prop:mult}
If $\tau$ is irreducible and nontrivial of dimension $d$, then
\[
m_\tau=-2d+\sum_{j=1}^rc_\tau(g_j),\qquad \dim\Hom_G(\tau,H^0(C,K_C))=\sum_{j=1}^r\age\tau(g_j^{-1})-d.
\]
\end{proposition}

\begin{proof}
Let $U=\PP^1-\{p_j\}$ and $\mathbb V$ the local system on $U$ attached to $\tau$. Then $m_\tau$ is the dimension of the parabolic cohomology $H^1(\PP^1,j_*\mathbb V)$, the image of $H^1_c(U,\mathbb V)\to H^1(U,\mathbb V)$. Since $\tau$ is irreducible and nontrivial, $H^0(U,\mathbb V)=0$, and $H^2(U,\mathbb V)=0$ because $U$ is an affine curve; so $\dim H^1(U,\mathbb V)=(r-2)d=\dim H^1_c(U,\mathbb V)$. The compactly supported dimension is also $(r-2)d$, since the Euler characteristic of $\mathbb V$ with compact supports is $\chi_c(U)d=(2-r)d$ and $H^0_c=H^2_c=0$, the latter by duality with $H^0(U,\mathbb V^\vee)=0$. The long exact sequence of the pair (the boundary circles of small punctured discs $\Delta_j^*$ around the $p_j$) begins
\[
0=H^0(U,\mathbb V)\to\textstyle\bigoplus_jH^0(\Delta_j^*,\mathbb V)\to H^1_c(U,\mathbb V)\to H^1(U,\mathbb V)\to\cdots,
\]
so the image of $H^1_c\to H^1$ has dimension $(r-2)d-\sum_j\dim H^0(\Delta_j^*,\mathbb V)$; and $\dim H^0(\Delta_j^*,\mathbb V)$ is the dimension $d-c_\tau(g_j)$ of the fixed space of $\tau(g_j)$. This gives the first formula. The second is the Chevalley--Weil formula \cite[Theorem~2.10]{FGP}.
\end{proof}

\subsection{The Prym--Tyurin factor}\label{subsec:pt}
Now let $\sigma$ be an irreducible representation of $G$ whose character field is $K$ and which is realizable over $K$, on $V_\sigma=K^d$. Here the character field is the field generated by the values of the character; since it is the imaginary quadratic field $K$, the character takes non-real values, so $\sigma\not\cong\bar\sigma$. The simple factor of $\Q[G]$ attached to $\sigma$ is $\Q[G]\varepsilon\cong\End_K(V_\sigma)\cong M_d(K)$, with centre $K$. Averaging gives a $G$-invariant positive definite $K$-Hermitian form $h_V$ on $V_\sigma$; the Rosati involution $g\mapsto g^{-1}$ acts on $\Q[G]\varepsilon$ as the $h_V$-adjoint and on its centre $K$ as complex conjugation.

\begin{definition}\label{def:pt}
Choose $v\in V_\sigma$ with $h_V(v,v)\ne0$, and let $e\in\Q[G]\varepsilon$ be the $h_V$-orthogonal projection onto the line $Kv$; it is a primitive idempotent fixed by the Rosati involution. For a cover $C$, the \emph{Prym--Tyurin factor} is
\[
A_\sigma(C)=\operatorname{im}\big(N_0e\colon J(C)\to J(C)\big),
\]
for an integer $N_0\ge1$ with $N_0e\in\Z[G]$, with the restricted polarization $E$ and the action $\iota$ of the centre $K$.
\end{definition}

Different choices of $v$ give $K$-isogenous factors. The $\sigma$-isotypic part of $J(C)$ is isogenous to $A_\sigma(C)^d$, one copy for each line of an orthogonal basis of $V_\sigma$. Prym--Tyurin varieties of this kind are classical \cite[Ch.~12]{BL}; Weil classes on them were first exploited for the quaternionic case by van Geemen and Verra \cite{vGV}.

\begin{lemma}\label{lem:weiltype}
$(A_\sigma(C),\iota,E)$ is a polarized abelian variety of dimension $m_\sigma$, and both embeddings of $K$ occur on its tangent space, with multiplicities the Hodge numbers $(p_\sigma,q_\sigma)$ of $M_\sigma$. If $m_\sigma=2n$ and $(p_\sigma,q_\sigma)=(n,n)$, it is a polarized abelian $2n$-fold of Weil type for $K$.
\end{lemma}

\begin{proof}
The $\sigma$-isotypic part of $H_1(J(C),\Q)$ is $V_\sigma\otimes_KM$ with $M=\Hom_{K[G]}(V_\sigma,H_1(J(C),\Q)_\varepsilon)$, and $H_1(A_\sigma,\Q)=Kv\otimes_KM\cong M$. After extending scalars along the two embeddings of $K$, $M$ becomes $M_\sigma$ and $M_{\bar\sigma}$. The Rosati involution of $E$ restricts on $\iota(K)$ to that of $\Q[G]$ on the centre, which is complex conjugation.
\end{proof}

\begin{remark}[$A_\sigma(C)$ is not of CM type]\label{rem:reflection}
In our application $(p_\sigma,q_\sigma)=(4,4)$ (Lemma~\ref{lem:g34mult}), and the family $A_\sigma(C)$ dominates the moduli space of split eightfolds (Section~\ref{sec:g34}). Since the Hodge group of a very general point of that moduli space is $SU_H$ \cite[Theorem~6.11]{vanGeemen}, so is that of $A_\sigma(C)$ for $C$ outside a countable union of proper analytic subsets of the component; in particular $A_\sigma(C)$ is not of CM type and not isogenous to a power of an elliptic curve. The equality $p_\sigma=q_\sigma$ is forced by the branch data: if every $g_j$ is conjugate to $g_j^{-1}$, then $\age\sigma(g_j)=\age\sigma(g_j^{-1})$, and Proposition~\ref{prop:mult}, applied to $\sigma$ and to $\bar\sigma$, gives $p_\sigma=q_\sigma$. This is consistent with Section~\ref{sec:brauer}: for real branch classes the line $\det M_\sigma$ is shown to be algebraic, so it must be of Hodge type.
\end{remark}

\begin{lemma}\label{lem:constancy}
On a connected component $S$ of $\cH$, the discriminant of the Hermitian form of $A_\sigma(C)$ does not depend on $C\in S$.
\end{lemma}

\begin{proof}
The spaces $H_1(A_\sigma(C),\Q)$ with $\iota$ and $E$ form a local system over $S$, since $e$ and $\iota$ act by flat endomorphisms and $E$ is flat; parallel transport is a $K$-isometry, and the discriminant is an isometry invariant.
\end{proof}

%% ------------------------------------------------------------------
\section{Algebraic Weil classes on the Prym--Tyurin factor}\label{sec:supply}
%% ------------------------------------------------------------------

Keep the notation of Section~\ref{sec:covers}, and assume $m_\sigma=2n$ and $(p_\sigma,q_\sigma)=(n,n)$. The idea of this section is the following. Schoen's theorem gives algebraic Weil classes on Prym varieties of cyclic covers. For a subgroup $H'\le G$ and a character $\chi$ of $H'$, the cyclic cover $C/\ker\chi\to C/H'$ is a quotient of $C$, so its Prym variety is a factor of $J(C)$. We look for an algebraic correspondence between a product of such Prym varieties and $A_\sigma(C)$ that carries the product of their Weil classes to a nonzero Weil class of $A_\sigma(C)$. Milne's theorem on Lefschetz classes (classes in the subalgebra generated by divisor classes) implies that a cohomology class on a product of factors of $J(C)$ is algebraic as soon as it is invariant under the centralizer of $G$, since that centralizer contains the relevant Lefschetz group. Invariance under the centralizer reduces, through its maximal abelian quotient, to an identity between characters, which is condition~(L) below.

\subsection{Cyclic suppliers and Schoen's theorem}
\begin{definition}\label{def:supplier}
A \emph{supplier} is a pair $(H',\chi)$ of a subgroup $H'\le G$ and a character $\chi\colon H'\to\C^\times$ of order $m\ge3$.
\end{definition}

For $C\in\cH$, the curve $C'=C/\ker\chi$ is a connected cyclic cover of $C/H'$ with group $H'/\ker\chi\cong\Z/m$. Its primitive part $P_\chi\subset J(C)$, the part on which $\Z/m$ acts through primitive characters, is cut out by a rational idempotent of $\Q[H']$, and the cyclotomic field $F_\chi=\Q(\mu_m)$ acts on it. For $t\in(\Z/m)^\times$ let $V_t\subset H^1(P_\chi,\C)$ be the eigenspace where $H'$ acts by $\chi^t$; then
\[
V_t=\bigoplus_\tau\Hom_{H'}(\chi^t,\tau|_{H'})\otimes M_\tau,
\]
all $V_t$ have the same dimension $h$, and the Weil classes $W_{F_\chi}(P_\chi)=\bigwedge^h_{F_\chi}H^1(P_\chi,\Q)$ satisfy $W_{F_\chi}(P_\chi)\otimes\C=\bigoplus_t\det V_t$.

The branch data of $C'\to C/H'$ are read off from the tuple. The points of $C/H'$ over $p_j$ correspond to the double cosets $H'\backslash G/\langle g_j\rangle$; over the double coset of $x$ the local monodromy is $xg_j^\ell x^{-1}\in H'$, with $\ell$ the least positive integer for which this lies in $H'$. Let $q$ be the genus of $C/H'$, $s$ the number of branch points of $C'\to C/H'$, and $\alpha\in(\Z/m-\{0\})^s$ their exponents with respect to a generator; then $h=2q-2+s$, which is the common dimension of the eigenspaces $V_t$ above (Chevalley--Weil for $C'\to C/H'$, \cite[Lemma~1.5]{Schoen88}). Following Schoen \cite[\S1]{Schoen88}, $\alpha$ is \emph{simple} if, after reordering, $\alpha_{2i}=-\alpha_{2i-1}$ for all $i$. A supplier is \emph{admissible} if $h$ is even and positive and $\alpha$ is simple (automatic when $s=0$).

\begin{lemma}\label{lem:classdatum}
For a supplier $(H',\chi)$ the numbers $q$, $s$, $h$, the exponents $\alpha$ up to $(\Z/m)^\times$ and order, and hence admissibility, depend only on $G$ and the classes $\mathcal C_j$.
\end{lemma}

\begin{proof}
Replacing $g_j$ by a conjugate $yg_jy^{-1}$ replaces the double coset of $x$ by that of $xy^{-1}$, with the same local monodromy; Riemann--Hurwitz for $C\to C/H'$ gives $q$.
\end{proof}

\begin{lemma}[Schoen]\label{lem:schoen}
If $(H',\chi)$ is admissible, then $W_{F_\chi}(P_\chi)$ is spanned by classes of algebraic cycles, for every $C\in\cH$.
\end{lemma}

\begin{proof}
We quote Schoen \cite[\S2]{Schoen88}. Let $C'$ be a smooth projective irreducible curve over $\C$ with a faithful action $\Z/m\hookrightarrow\operatorname{Aut}(C')$, $m>1$, whose invariant $h$ is even, so that the quotient map has an even number $2r$ of branch points, with exponents $1\le\alpha_j\le m-1$. \emph{Theorem~2.0}: if the $(\Z/m)^\times$-orbit of $(\alpha_1,\dots,\alpha_{2r})$ is simple, then the relevant subspace of $H^h(C'^h)$ is generated by fundamental classes of algebraic cycles. \emph{Corollary~3.1}, for $m>2$: under the same hypotheses, the corresponding subspace of $H^h(B,\Q)$ is generated by classes of algebraic cycles of codimension $h/2$, where $B\subseteq\operatorname{Alb}(C')$ is the image of $\prod_{1\ne Z'\le\Z/m}\sum_{t\in Z'}(\mathrm{Id}-t)$, the primitive part. Neither statement asks the exponents to be prime to $m$: Schoen's \S1 treats a branch point with $d=\gcd(\alpha_j,m)>1$ through the local model $y^{m/d}=w^{\alpha_j/d}$. Simplicity is automatic for $r\le1$, and the unramified case $r=0$ is included, with its own proof \cite[\S2]{Schoen88}.

Apply this to $C'=C/\ker\chi\to C/H'$, which is smooth, projective and irreducible as a quotient of $C$, with $H'/\ker\chi\cong\Z/m$ acting faithfully. The subspace of Corollary~3.1 is $W_{F_\chi}$ of the primitive part, which is isogenous to $P_\chi$ compatibly with $F_\chi$. Admissibility does not depend on $C$ by Lemma~\ref{lem:classdatum}. When every exponent is $m/2$ and $4\mid m$, one step of Schoen's proof is replaced by Lemma~\ref{lem:schoenR} below.
\end{proof}

For the group $G_{34}$ of Section~\ref{sec:g34} every local monodromy is an involution, so every branch point of every supplier cover has exponent $m/2$, and $4\mid m$ occurs: three of the $84$ suppliers of the explicit identity of Remark~\ref{rem:witness} have $m=12$ and $s>0$ (rows $42$--$44$ of the deposited rebuild log). For such covers one step of Schoen's proof needs a replacement, which we now give.

Recall Schoen's construction \cite[\S2, pp.~11--19]{Schoen88}, for $C'\to X$ as in Lemma~\ref{lem:schoen} with branch points ordered so that $\alpha_{2i}=-\alpha_{2i-1}$, $1\le i\le r$. Let $W_0=C'^h/G_0$, where $G_0\subset(\Z/m)^h\rtimes S_h$ is the kernel of the sum map extended by $S_h$; the residual group $\Z/m$ acts on $W_0$, and $\beta_0\colon W_0\to S^hX$ is a $\Z/m$-cover branched along the divisors $S_j=b_j+S^{h-1}X$. Schoen chooses a linear subspace $P_0\subset S^hX$ of dimension $h/2$ (his Lemma~2.3) and, for $i=1,\dots,r$, blows up $(S^hX)_{i-1}$ along $K_{i-1}=S_{2i-1,i-1}\cap S_{2i,i-1}$, the intersection of the strict transforms of $S_{2i-1}$ and $S_{2i}$; the exceptional divisor $E_i$ is a $\PP^1$-bundle over $K_{i-1}$. Let $W_i$ be the normalization of the pulled-back cover, $\beta_i\colon W_i\to(S^hX)_i$, and $F_i=\beta_i^{-1}(E_i)$. After $r$ steps the strict transform of $P_0$ avoids the branch locus, its inverse image splits into $m$ components, and these give the cycles of Theorem~2.0 (his Lemmas~2.4--2.6). To see that their classes on $W_0$ are nonzero, his Lemma~2.8 shows that no primitive character of $\Z/m$ occurs in $H^*(W_i,\C)/H^*(W_{i-1},\C)$, and this reduces to the statement that no primitive character occurs in $H^*(F_i,\C)$ \cite[p.~18]{Schoen88}. Schoen proves the latter (Sublemma~2.9, p.~19) by choosing a prime $p\mid m$ such that the $\Z/p$-quotient of $C'\to X$ is ramified over $b_{2i-1}$. When every exponent equals $m/2$ and $4\mid m$, no such prime exists.

\begin{lemma}[Schoen's Theorem~2.0 for exponents $m/2$]\label{lem:schoenR}
In the situation of Schoen's Theorem~2.0, for every $i\le r$ and every primitive character $\chi$ of $\Z/m$, the $\chi$-part of $H^*(F_i,\C)$ vanishes. Hence Theorem~2.0 and Corollary~3.1 of \cite{Schoen88} hold under their stated hypotheses, without the choice of a prime in Sublemma~2.9; in particular they hold when every exponent is $m/2$ and $4\mid m$.
\end{lemma}

\begin{proof}
Since $\beta_i$ is finite, $H^*(F_i,\C)_\chi=H^*(E_i,(\beta_{i*}\C)_\chi|_{E_i})$. The stalk of $\beta_{i*}\C$ at a point is the permutation module on the points above it, and its $\chi$-part vanishes at every point whose stabilizer contains nontrivial inertia, because a primitive character is faithful. By Schoen's Lemma~2.5, $E_i$ is not a branch divisor of $\beta_i$: its multiplicity in the branch locus is $\alpha_{2i-1}+\alpha_{2i}\equiv0\pmod m$. So on $E_i$ the sheaf $(\beta_{i*}\C)_\chi$ is $j_!\mathcal L$, where $\mathcal L$ is a rank-one local system on the complement in $E_i$ of the strict transforms $S_{j,i}$, with monodromy $\chi(\alpha_j)\ne1$ around $S_{j,i}$. The strict transforms $S_{2i-1,i}$ and $S_{2i,i}$ meet $E_i$ in two disjoint sections of $p\colon E_i\to K_{i-1}$, and every other $S_{j,i}$ meets $E_i$ in a union of whole fibres. Such a fibre lies in the closed set $S_{j,i}$, so the restriction of $j_!\mathcal L$ to it is the zero sheaf. On any other fibre $\PP^1$ it is the extension by zero of a rank-one local system on $\C^\times=\PP^1-\{0,\infty\}$ with monodromy $\chi(\alpha_{2i-1})\ne1$, whose compactly supported cohomology vanishes: $H^0_c=0$, $H^2_c$ is dual to the invariants of the dual system, which are $0$, and the Euler characteristic is $0$. Since $p$ is proper and $j_!\mathcal L$ is constructible, proper base change identifies the stalk of $R^bp_*(j_!\mathcal L)$ at $k\in K_{i-1}$ with $H^b(p^{-1}(k),j_!\mathcal L|_{p^{-1}(k)})$ \cite[Theorem~11.3.6]{Schapira}, which we have just shown to vanish. Hence $Rp_*(j_!\mathcal L)=0$ and $H^*(E_i,j_!\mathcal L)=H^*(K_{i-1},Rp_*j_!\mathcal L)=0$. This uses only $\alpha_{2i-1}\not\equiv0\pmod m$. The rest of Schoen's proof is unchanged.
\end{proof}

\subsection{The character condition}
Let $U$ be the centralizer of $G$ in the symplectic group of $(H^1(C,\Q),E)$, an algebraic group over $\Q$. We only need its complex points. Recall that an irreducible complex representation $\tau$ of a finite group is \emph{orthogonal} or \emph{quaternionic} if it is self-dual and carries a nonzero invariant symmetric, respectively alternating, bilinear form (Frobenius--Schur indicator $+1$ or $-1$), and that otherwise $\tau\not\cong\bar\tau$ \cite[\S13.2]{SerreLin}. In $H^1(C,\C)=\bigoplus_\tau\tau\otimes M_\tau$ the $G$-invariant alternating cup-product form pairs $\tau\otimes M_\tau$ with $\bar\tau\otimes M_{\bar\tau}$ and is zero between other summands. By Schur's lemma its restriction is the tensor product of the invariant form on $\tau$, unique up to scalar, with a nondegenerate form on the multiplicity spaces: a duality between $M_\tau$ and $M_{\bar\tau}$ if $\tau\not\cong\bar\tau$; a symmetric form on $M_\tau$ if $\tau$ is quaternionic; an alternating one if $\tau$ is orthogonal. An automorphism commuting with $G$ acts through $\prod_\tau GL(M_\tau)$, and preserving the form means preserving these forms. Hence $U(\C)$ is the product of the groups $GL(M_\tau)$ over the pairs $\{\tau,\bar\tau\}$ of non-self-dual irreducible representations with $m_\tau>0$ (acting on $M_\tau$ and contragrediently on $M_{\bar\tau}$), of $O(M_\tau)$ over the quaternionic $\tau$ with $m_\tau>0$, and of $Sp(M_\tau)$ over the orthogonal $\tau$ with $m_\tau>0$; representations with $m_\tau=0$ do not occur in $H^1(C)$ and play no role. Let $T=U(\C)/[U(\C),U(\C)]$ be its maximal abelian quotient. Its character group is
\[
X^*(T)=\Z^{\rm pairs}\oplus(\Z/2)^{\rm quat},
\]
where ``pairs'' and ``quat'' run over the pairs and the quaternionic $\tau$ with $m_\tau>0$: one free coordinate for each pair (the determinant on $M_\tau$, for a chosen member $\tau$ of the pair; for the pair of $\sigma$ we choose $\sigma$), and one coordinate of order two (the determinant of $O(M_\tau)$) for each quaternionic $\tau$. Each line $\det V_t$ above is an eigenline of $U$; its character is
\[
a_t(\chi)=\big(\langle\tau|_{H'},\chi^t\rangle-\langle\bar\tau|_{H'},\chi^t\rangle\big)_{\rm pairs}\oplus\big(\langle\tau|_{H'},\chi^t\rangle\bmod2\big)_{\rm quat},
\]
and we put $a(\chi)=a_1(\chi)$. The line $\det M_\sigma\subset H^{2n}(A_\sigma(C),\C)$, one of the two lines of $W_K(A_\sigma)\otimes\C$, has character $e_\sigma$: coordinate $+1$ at the pair of $\sigma$ and $0$ elsewhere. Complex conjugation exchanges $\chi$ and $\chi^{-1}$, so $a_{-1}(\chi)=-a(\chi)$ on the free coordinates and $a_{-1}(\chi)=a(\chi)$ on the others.

\begin{definition}\label{def:lattice}
Let $L(G;\mathcal C_1,\dots,\mathcal C_r)\subseteq X^*(T)$ be the subgroup generated by the $a(\chi)$ of the admissible suppliers. \emph{Condition (L)} is that $e_\sigma\in L$.
\end{definition}

By Lemma~\ref{lem:classdatum}, condition (L) depends only on $G$, $\sigma$ and the classes $\mathcal C_j$.

\begin{proposition}\label{prop:supply}
Let $C\in\cH$ and let $X=\prod_kX_k$ be a finite product of abelian subvarieties $X_k=\operatorname{im}(N_ke_k)\subseteq J(C)$, with $e_k$ idempotents of $\Q[G]$. Let $\ell\subset H^p(X,\C)$ be an eigenline of $U$ whose character lies in $L$. Then $\ell$ lies in the complex span of classes of algebraic cycles on $X$.
\end{proposition}

\begin{proof}
\emph{Spaces, classes and degrees.} Write the character of $\ell$ as $\xi=\sum_in_ia(\chi_i)$ with admissible $(H'_i,\chi_i)$ and $n_i\ne0$. Put $t_i=\operatorname{sign}n_i$, $P_i=P_{\chi_i}$, $m_i=\operatorname{ord}\chi_i$, $h_i=h(\chi_i)$, so $\dim P_i=\varphi(m_i)h_i/2$, and let $V_{i,t}$, $t\in(\Z/m_i)^\times$, be the eigenspaces of $H^1(P_i,\C)$, each of dimension $h_i$. Let
\[
B=\textstyle\prod_iP_i^{|n_i|},\qquad d_B=\dim B=\sum_i|n_i|\,\varphi(m_i)h_i/2 .
\]
By Lemma~\ref{lem:schoen} choose a nonzero rational algebraic class $y_i\in W_{F_{\chi_i}}(P_i)\subset H^{h_i}(P_i,\Q)$, and put $\mathbf y=\boxtimes_iy_i^{\boxtimes|n_i|}\in H^{h_B}(B,\Q)$, $h_B=\sum_i|n_i|h_i$, an algebraic class. Choose generators $z_{i,t}$ of the lines $\det V_{i,t}$ and put $Z_i=\prod_{t\ne t_i}z_{i,t}\in H^{(\varphi(m_i)-1)h_i}(P_i,\C)$. Let $x$ generate $\ell\subset H^p(X,\C)$ and
\[
\Gamma=\Big(\textstyle\boxtimes_iZ_i^{\boxtimes|n_i|}\Big)\boxtimes x\in H^{2d_B-h_B+p}(B\times X,\C).
\]
For a class $\Gamma$ on $B\times X$, the associated correspondence is $\Gamma_*\colon H^k(B)\to H^{k+\deg\Gamma-2d_B}(X)$, $\Gamma_*(\beta)=\operatorname{pr}_{2*}(\operatorname{pr}_1^*\beta\cup\Gamma)$. If $\Gamma$ is a complex combination of algebraic classes and $\beta$ is algebraic, $\Gamma_*(\beta)$ is a complex combination of algebraic classes. Here $\Gamma_*(\mathbf y)\in H^p(X,\C)$.

\emph{Step 1: $\Gamma$ is $U$-invariant.} The product $\prod_tz_{i,t}$ spans the top cohomology of $P_i$, on which $U$ acts trivially, since $U$ preserves the polarization; so $Z_i$ has character $-a(\chi_i^{t_i})$, which is $-t_ia(\chi_i)$ on the free coordinates and $a(\chi_i)$ on the others. Hence $\Gamma$ has character $\xi-\sum_in_ia(\chi_i)=0$ on the free coordinates, and $\xi-\sum_i|n_i|a(\chi_i)\equiv0$ on the coordinates of order two.

\emph{Step 2: $\Gamma$ is algebraic.} Let $R$ be the number of factors of $B\times X$, namely $\sum_i|n_i|$ copies of the $P_i$ and the factors $X_k$. Each factor $Y'$ is $\operatorname{im}(N_0e')$ for an idempotent $e'\in\Q[G]$ (for $P_i$, the rational idempotent of $\Q[H'_i]$ cutting out $P_{\chi_i}$), with one integer $N_0$ for all factors such that every $N_0e'\in\Z[G]$. Let $\iota_{Y'}\colon Y'\to J(C)$ be the inclusion and $\pi_{Y'}\colon J(C)\to Y'$ the map $N_0e'$ with its target restricted to its image; then $\pi_{Y'}\circ\iota_{Y'}=[N_0]$, because $e'$ acts as the identity on its image. Taking products,
\[
B\times X\ \xrightarrow{\ \iota=\prod\iota_{Y'}\ }\ J(C)^R\ \xrightarrow{\ \pi=\prod\pi_{Y'}\ }\ B\times X,\qquad \pi\circ\iota=[N_0].
\]
The group $U$ commutes with $\Q[G]$, so it preserves each $H^1(Y')=e'H^1(J(C))$, and $\iota^*$ and $\pi^*$ are $U$-equivariant, with $U$ acting diagonally on $J(C)^R$. So $\pi^*\Gamma\in H^*(J(C)^R,\C)$ is $U$-invariant. Milne defines $S(A)$, for an abelian variety $A$ and the symplectic form $e_D$ of an ample divisor, as the largest algebraic subgroup of $Sp(e_D)$ whose elements commute with $\End^0(A)$ \cite[\S1]{Milne}; it acts diagonally on $H^*(A^r)$, and his Theorem~3.2 states: \emph{for any abelian variety $A$ and integer $r\ge0$, the algebra $H^*(A^r)^{S(A)}$ is generated by divisor classes}; by his Lemma~3.1 this holds with complex coefficients. For $A=J(C)$ with the canonical polarization, $S(J(C))$ commutes with $\Q[G]\subseteq\End^0(J(C))$, so $S(J(C))\subseteq U$, and $\pi^*\Gamma$ is $S(J(C))$-invariant. By Theorem~3.2 with $r=R$, $\pi^*\Gamma$ is a complex combination of products of divisor classes on $J(C)^R$. Since $[N_0]^*$ acts on $H^k$ of an abelian variety as multiplication by $N_0^k$, we get $\iota^*\pi^*\Gamma=[N_0]^*\Gamma=N_0^{\deg\Gamma}\,\Gamma$, so $\Gamma$ is a complex combination of algebraic classes.

\emph{Step 3: the contraction is nonzero.} Write $y_i=\sum_tc_{i,t}v_{i,t}$ with $v_{i,t}\in\det V_{i,t}$ normalized by $\int_{P_i}v_{i,t_i}\cup Z_i=1$. For $t\ne t_i$ the product $v_{i,t}\cup Z_i$ contains $\det V_{i,t}$ twice and vanishes. By the K\"unneth formula and the projection formula,
\[
\Gamma_*(\mathbf y)=\pm\textstyle\prod_ic_{i,t_i}^{|n_i|}\,x ,
\]
the sign coming from reordering graded factors. Each $c_{i,t}$ is nonzero, because $W_{F_{\chi_i}}(P_i)$ is a one-dimensional $F_{\chi_i}$-vector space and a nonzero vector in it has a nonzero component under every embedding of $F_{\chi_i}$. So $x$ lies in the complex span of algebraic classes. Neither the complements $Z_i$ nor $\Gamma$'s factors are claimed to be algebraic individually.
\end{proof}

\begin{corollary}\label{cor:everymember}
If condition (L) holds, then $W_K(A_\sigma(C))$ is spanned by classes of algebraic cycles for every $C\in\cH$.
\end{corollary}

\begin{proof}
$W_K(A_\sigma)\otimes\C=\det M_\sigma\oplus\det M_{\bar\sigma}$, two eigenlines of $U$ with characters $\pm e_\sigma\in L$. Proposition~\ref{prop:supply} puts both in the complex span of algebraic classes. If $\mathcal Z$ is the rational span of algebraic classes and $W\otimes\C\subseteq\mathcal Z\otimes\C$ for a rational subspace $W$, then $W\subseteq\mathcal Z$.
\end{proof}


%% ------------------------------------------------------------------
\section{Real branch classes and Brauer's induction theorem}\label{sec:brauer}
%% ------------------------------------------------------------------

Call the classes $\mathcal C_j$ \emph{real} if every $g_j$ is conjugate in $G$ to $g_j^{-1}$; involutions are real.

\begin{theorem}\label{thm:brauer}
If the classes $\mathcal C_1,\dots,\mathcal C_r$ are real, condition (L) holds.
\end{theorem}

\begin{proof}
\emph{A homomorphism.} Define $\mathsf a\colon R_\C(G)\to X^*(T)$ on the representation ring by
\[
\mathsf a(V)=\big(\operatorname{mult}_\tau V-\operatorname{mult}_{\bar\tau}V\big)_{\rm pairs}\oplus\big(\operatorname{mult}_\eta V\bmod2\big)_{\rm quat}.
\]
It is additive, $\mathsf a(\sigma)=e_\sigma$, and by Frobenius reciprocity $\mathsf a(\Ind_{H'}^G\chi)=a(\chi)$ for every subgroup $H'$ and linear character $\chi$ of $H'$.

\emph{Brauer induction.} By Brauer's theorem \cite[\S10.5, Theorem~20]{SerreLin}, $\sigma=\sum_in_i\Ind_{B_i}^G\lambda_i$ in $R_\C(G)$ with $n_i\in\Z$, $B_i\le G$ elementary (a product of a cyclic group and a $p$-group for some prime $p$) and $\lambda_i$ linear. Hence $e_\sigma=\sum_in_ia(\lambda_i)$.

\emph{Terms that vanish.} If $\lambda_i$ has order $1$ or $2$ it is real, so the free coordinates of $a(\lambda_i)$ vanish; for quaternionic $\eta$ the coordinate of order two is $\langle\eta|_{B_i},\lambda_i\rangle\bmod2$, and it vanishes. Indeed, let $\langle\ ,\ \rangle_\eta$ be the nondegenerate $G$-invariant alternating form on $\eta$, and decompose $\eta|_{B_i}=\bigoplus_\rho\eta_\rho$ into isotypic parts over the irreducible representations $\rho$ of $B_i$. An invariant pairing between $\eta_\rho$ and $\eta_{\rho'}$ is a $B_i$-map $\eta_{\rho'}\to\eta_\rho^\vee$, which vanishes unless $\rho'\cong\rho^\vee$; so the form pairs $\eta_\rho$ only with $\eta_{\rho^\vee}$. Since $\lambda_i^\vee=\lambda_i^{-1}=\lambda_i$, the space $\eta_{\lambda_i}$ is orthogonal to all the other isotypic parts, so the form is nondegenerate on it; an alternating nondegenerate form lives on an even-dimensional space, so $\langle\eta|_{B_i},\lambda_i\rangle=\dim\eta_{\lambda_i}$ is even. If $\lambda_i$ has order at least $3$ and $h(\lambda_i)=0$, then every $\tau$ with $m_\tau>0$ has $\langle\tau|_{B_i},\lambda_i\rangle=0$ and $a(\lambda_i)=0$. So $e_\sigma$ is an integer combination of $a(\lambda_i)$ over suppliers with $h>0$.

\emph{These suppliers are admissible.} Let $(H',\chi)$ be a supplier with $h>0$, and $w\in G$ with $wg_jw^{-1}=g_j^{-1}$. Since $w$ normalizes $D=\langle g_j\rangle$, the map $H'xD\mapsto H'xwD$ is a bijection of $H'\backslash G/D$, and $(xw)g_j^k(xw)^{-1}=xg_j^{-k}x^{-1}$; so it preserves the integer $\ell$ above and inverts the local monodromy. Hence the exponents of $C'\to C/H'$ over $p_j$ form a multiset stable under $\alpha\mapsto-\alpha$. Pair each exponent $\alpha\ne-\alpha$ with an occurrence of $-\alpha$; the unpaired ones equal $m/2$. Their number $k$ is even, because the exponents sum to $0$ modulo $m$ and the pairs contribute $0$, so $k\cdot m/2\equiv0\pmod m$. Hence $\alpha$ is simple, $s$ is even, and $h=2q-2+s$ is even.
\end{proof}

%% ------------------------------------------------------------------
\section{Dominance and spreading}\label{sec:spread}
%% ------------------------------------------------------------------

Fix a component $S$ of $\cH$ and let $\fd(S)$ be the discriminant of $A_\sigma(C)$, $C\in S$ (Lemma~\ref{lem:constancy}). Let $\cM=\cM(K,n,\fd(S))$ be the moduli space of Section~\ref{subsec:moduli}. Locally on $S$ the factors $A_\sigma(C)$ give a holomorphic period map to $\cHn$, well defined up to $\Gamma(\ell)$; we say that $S$ is \emph{dominant} if this map has rank $n^2$ at some point, hence on a dense open subset.

\begin{lemma}\label{lem:locus}
Let $\mathcal W=\bigwedge^{2n}_KR^1f_*\Q\subset R^{2n}f_*\Q$ on $\cM$. The set $\cM^{\rm alg}$ of points $t$ at which $\mathcal W_t$ is spanned by classes of algebraic cycles is a countable union of closed algebraic subsets of $\cM$.
\end{lemma}

\begin{proof}
This is the standard argument with relative Hilbert schemes \cite[\S3.3.2, pp.~90--94, and the proofs of Theorems~7.19 and 8.26]{VoisinII}; we give the details. Fix a relatively ample line bundle on the smooth projective family $f\colon\mathcal A\to\cM$. For each Hilbert polynomial, the relative Hilbert scheme of $n$-dimensional closed subschemes of the fibres is projective over $\cM$; there are countably many Hilbert polynomials, and each relative Hilbert scheme has finitely many irreducible components \cite[pp.~90--91]{VoisinII}. A rational algebraic class in $H^{2n}(\mathcal A_t,\Q)$ is a rational combination of classes of $n$-dimensional subschemes of $\mathcal A_t$, so it suffices to consider finite tuples of such subschemes.

Fix a finite tuple $\mathcal H_1,\dots,\mathcal H_k$ of components, let $T$ be an irreducible component of $\mathcal H_1\times_\cM\dots\times_\cM\mathcal H_k$, let $\widetilde T\to T$ be a resolution of singularities, and let $g\colon\widetilde T\to\cM$ be the composite. Then $\widetilde T$ is smooth and connected, and $g$ is proper, so $g(\widetilde T)$ is a closed algebraic subset of $\cM$. The universal subschemes pull back to subschemes $\mathcal Z_l\subset\mathcal A\times_\cM\widetilde T$, flat over $\widetilde T$, and the classes $[\mathcal Z_{l,u}]\in H^{2n}(\mathcal A_{g(u)},\Q)$, $u\in\widetilde T$, are global flat sections of the local system $g^*R^{2n}f_*\Q$ on $\widetilde T$: the fibre classes of a flat family of subschemes of a smooth projective family are locally constant, as used in \cite[proof of Theorem~3.33, p.~94]{VoisinII}. Let $\mathcal F_u$ be the span of $[\mathcal Z_{1,u}],\dots,[\mathcal Z_{k,u}]$. Parallel transport along a path from $u$ to $u'$ is an isomorphism $(g^*R^{2n}f_*\Q)_u\to(g^*R^{2n}f_*\Q)_{u'}$ carrying each of these global sections to itself, so it carries $\mathcal F_u$ onto $\mathcal F_{u'}$ and $g^*\mathcal W_u$ onto $g^*\mathcal W_{u'}$; in particular $\dim\mathcal F_u$ does not jump. Hence the set of $u$ with $g^*\mathcal W_u\subseteq\mathcal F_u$ is either all of $\widetilde T$ or empty. Call $T$ \emph{good} in the first case; then $\mathcal W_t$ is algebraic for every $t\in g(\widetilde T)$.

Conversely, if $\mathcal W_t$ is spanned by algebraic classes, choose subschemes $Z_1,\dots,Z_k$ of $\mathcal A_t$ whose classes span a space containing $\mathcal W_t$; the point $(Z_1,\dots,Z_k)$ lies on some component $T$, which is then good, and $t\in g(\widetilde T)$. So $\cM^{\rm alg}$ is the union of the closed sets $g(\widetilde T)$ over the good $T$, of which there are countably many.
\end{proof}

\begin{theorem}\label{thm:transfer}
Assume condition (L), and let $S$ be a dominant component of $\cH$. Then the Weil classes of every polarized abelian $2n$-fold of Weil type for $K$ with discriminant $\fd(S)$ are algebraic.
\end{theorem}

\begin{proof}
Choose $C_1\in S$ where the period map has rank $n^2$ and a small polydisc $D\ni C_1$ over which the covers form an analytic family. Over the simply connected $D$ the spaces $H_1(A_\sigma(C),\Q)$ with $\iota$ and $H$ form a trivial flat bundle of Hermitian spaces of discriminant $\fd(S)$, so Landherr's theorem gives a $K$-isometry $\mu\colon H_1(A_\sigma(C_1),\Q)\to(V,H)$, which we transport over $D$. Sending $C$ to the complex structure of $H_1(A_\sigma(C),\R)$, transported by $\mu$, gives the period map $\Phi\colon D\to\cHn$, a submersion at $C_1$. Shrinking $D$, $\Phi(D)$ contains a ball $\mathcal B$, whose image in $\cM$ is open, since $\Gamma(\ell)$ acts properly and freely.

\emph{The fibres over the ball.} Let $C\in D$ with $z=\Phi(C)\in\mathcal B$, and let $A_z$ be the fibre of $\mathcal A$ over the image of $z$: $H_1(A_z,\Z)=\Lambda$, with complex structure $z$ and polarization $\Tr_{K/\Q}(H/2\vartheta)$ scaled as in Subsection~\ref{subsec:moduli}. The lattice $\mu(H_1(A_\sigma(C),\Z))$ is commensurable with $\Lambda$; if $N\mu(H_1(A_\sigma(C),\Z))\subseteq\Lambda$ for an integer $N\ge1$, then $N\mu$ is a $K$-linear map of lattices compatible with the complex structures, so it is a $K$-linear isogeny $A_\sigma(C)\to A_z$. Since $\mu$ is an isometry of Hermitian forms, it pulls the polarization of $A_z$ back to a positive rational multiple of that of $A_\sigma(C)$; neither the Weil plane nor the class of the discriminant changes. By Corollary~\ref{cor:everymember} and Lemma~\ref{lem:isogeny}, the Weil classes of $A_z$ are algebraic. So $\cM^{\rm alg}$ contains a nonempty open set. A nonempty open subset of $\cM$ is a Baire space, so by Lemma~\ref{lem:locus} one of the closed sets has nonempty interior, and since $\cM$ is irreducible it is all of $\cM$. Every polarized $2n$-fold of Weil type for $K$ with discriminant $\fd(S)$ is $K$-isogenous to a fibre of $\mathcal A$ (Lemma~\ref{lem:isogeny}).
\end{proof}

%% ------------------------------------------------------------------
\section{The Hermitian form as a parabolic cup product}\label{sec:herm}
%% ------------------------------------------------------------------

To compute with the family $A_\sigma(C)$ we need its first cohomology and Hermitian form in terms of the monodromy tuple alone.

Let $\pi\colon C\to\PP^1$ be a cover in $\cH$, $U_C=\PP^1-\{p_j\}$, and $\mathcal R=\pi_*\Q|_{U_C}$, a local system with fibre $\Q[G]$. Put $L'=\Hom_G(V_\sigma,\mathcal R\otimes K)$, a local system of $K$-vector spaces of rank $d$ whose monodromy is $\sigma$ composed with that of $\pi$. The fibre pairing of $\mathcal R$ (multiplication followed by the sum over the fibre), made conjugate-linear in the first variable, restricts on $V_\sigma\otimes L'$ to $h_V\otimes h_{L'}$ for a flat positive definite Hermitian form $h_{L'}$ on $L'$, unique up to a positive rational factor. Let
\[
\Vp=H^1_{\rm par}(U_C,L'),\qquad \mathcal S(a,b)=\int_{U_C}h_{L'}(a\cup b),\qquad h=\vartheta\mathcal S,
\]
the parabolic cohomology (the image of compactly supported cohomology), a $K$-space of dimension $m_\sigma$, with its skew-Hermitian cup product $\mathcal S$ and the Hermitian form $h$.

\begin{proposition}\label{prop:herm}
There is a $K$-linear or $K$-antilinear isomorphism $H_1(A_\sigma(C),\Q)\cong\Vp$ under which the Hermitian form $H$ of $A_\sigma(C)$ is a nonzero rational multiple of $h$, or of $\bar h$. In particular $H$ and $h$ have the same discriminant, and one is split if and only if the other is.
\end{proposition}

\begin{proof}
\emph{Step 1: the skew-Hermitian lift.} On $H_1(A,\Q)$ the form $T=H/2\vartheta$ is the unique $K$-sesquilinear skew-Hermitian form with $\Tr_{K/\Q}T=E$.

\emph{Step 2: the multiplicity space.} Let $N'=H_1(J(C),\Q)_\varepsilon$ and $M=\Hom_{K[G]}(V_\sigma,N')$. Evaluation $V_\sigma\otimes_KM\to N'$ is a $G$-equivariant isomorphism, because $\End_{K[G]}(V_\sigma)=K$. A $G$-invariant sesquilinear form on $N'$ pulls back to $h_V\otimes\psi$ for a unique sesquilinear form $\psi$ on $M$, since the $G$-invariant sesquilinear forms on $V_\sigma$ are the multiples of $h_V$. Applied to $T$, which is $G$-invariant by the uniqueness in Step~1, this gives $T=h_V\otimes\psi$, and $T|_{A_\sigma}=h_V(v,v)\,\psi$ under $m\mapsto v\otimes m$, because $e$ is Rosati-fixed.

\emph{Step 3: the cup product.} Since $\pi$ is finite and $C$ is normal, $H^1(C,K)=H^1_{\rm par}(U_C,\mathcal R\otimes K)$, with the cup product given by the fibre pairing. The functor $\Hom_{K[G]}(V_\sigma,-)$ is exact and commutes with these operations, so it identifies the $\sigma$-part of the skew-Hermitian lift of the cup product with $h_V\otimes\mathcal S$.

\emph{Step 4: comparison.} Poincar\'e duality on $C$ and the identification of $H_1$ with $H^1$ by $E$ are $G$-equivariant and carry $T$ to $\pm$ the lift of Step~3. They carry the $\sigma$-component either to the $\sigma$-component or to the $\bar\sigma$-component, depending on sign conventions that we do not need to fix; accordingly the induced isomorphism $M\cong\Vp$ is $K$-linear or $K$-antilinear, and $\psi=c\,\mathcal S$ or $\psi=c\,\bar{\mathcal S}$ with $c\in K^\times$. Both forms are skew-Hermitian, so $c=\bar c$ is rational. In either case the discriminants agree: scaling a Hermitian form of even dimension by a rational number multiplies its determinant by a norm, and $\det\bar h=\det h$.
\end{proof}

The form $\mathcal S$ is computable from the tuple. Let $W$ be a local system of $K$-vector spaces on $\PP^1$ minus $r$ points with a flat Hermitian form $h_W$ and unitary monodromy $T_1,\dots,T_r$ along standard loops with $\ell_1\cdots\ell_r=1$, generating a finite group. Put $Q_i=T_1\cdots T_{i-1}$. A parabolic cocycle is given by vectors $a_1,\dots,a_r$ with $\sum_iQ_i(1-T_i)a_i=0$; its value on $\ell_i$ is $c_i=(1-T_i)a_i$. Put $F_i=\sum_{j<i}Q_jc_j$.

\begin{lemma}\label{lem:cocycle}
For parabolic cocycles $a,b$, the cup product of their classes with the fibre pairing $h_W$, on the fundamental class oriented so that the loops run clockwise, is
\[
\mathcal S(a,b)=\sum_{i=1}^rh_W\big(Q_i^{-1}F^a_i+a_i,\ c^b_i\big).
\]
It does not depend on the choice of the vectors $a_i$, which are determined by the cocycle only modulo $\ker(1-T_i)$; it vanishes on coboundaries and is skew-Hermitian.
\end{lemma}

\begin{figure}[ht]
\centering
\begin{tikzpicture}[scale=1.2,>=stealth]
\foreach \a/\n in {90/0,0/1,-90/2,180/3} {
  \draw[->] (0,0) -- (\a:2.9);
  \node at (\a:3.3) {$e_{\n}$};
}
\foreach \a/\n in {45/1,-45/2,-135/3,135/4} {
  \fill (\a:1.4) circle (2pt);
  \node at ($(\a:1.4)+(\a+90:0.32)$) {$p_{\n}$};
  \draw[->] (0,0) -- (\a:1.33);
  \node at ($(\a:0.75)+(\a-90:0.2)$) {$f_{\n}$};
  \draw[->,dashed] (\a:1.47) -- (\a:2.9);
  \node at ($(\a:2.25)+(\a-90:0.22)$) {$g_{\n}$};
}
\fill (0,0) circle (2pt);
\node at (-67.5:0.42) {$b_0$};
\node at (3.6,-2.6) {$b_\infty=\infty$};
\end{tikzpicture}
\caption{The $\Delta$-complex of Lemma~\ref{lem:cocycle} for $r=4$, drawn in the plane $\PP^1-\{\infty\}$. The edges $e_i$ and $g_i$ end at $b_\infty=\infty$. The sector between $e_{i-1}$ and $e_i$ contains $p_i$ and is cut by the path $f_i\cup g_i$ into the two triangles at $p_i$. The loop $\ell_i$, out along $e_{i-1}$ and back along $e_i$, runs clockwise around $p_i$.}\label{fig:delta}
\end{figure}

\begin{proof}
We use the following $\Delta$-complex structure on $\PP^1$ (Figure~\ref{fig:delta}). Its vertices are a base point $b_0$, a second point $b_\infty$, and the punctures $p_1,\dots,p_r$. Its edges are $e_0,\dots,e_{r-1}$ from $b_0$ to $b_\infty$, cutting the sphere into $r$ sectors, with the sector between $e_{i-1}$ and $e_i$ containing $p_i$ (indices of $e$ modulo $r$, so $e_r=e_0$); and, for each $i$, an edge $f_i$ from $b_0$ to $p_i$ and an edge $g_i$ from $p_i$ to $b_\infty$ inside that sector. Its faces are the $2r$ triangles with vertices $(b_0,p_i,b_\infty)$ in this order: for each $i$ one triangle with edges $f_i,g_i,e_{i-1}$ and one with edges $f_i,g_i,e_i$. (The Euler characteristic is $(r+2)-3r+2r=2$.) The loop $\ell_i$ goes out along $e_{i-1}$ and back along $e_i$, and we orient $\PP^1$ so that it runs clockwise around $p_i$; then $\ell_1\cdots\ell_r$ is out along $e_0$ and back along $e_r=e_0$, which is trivial.

The monodromy group is finite, so there is a finite Galois cover branched over the punctures on which $W$ becomes trivial; lift the complex to it. There the cup product of $1$-cochains followed by $h_W$, evaluated on the fundamental class, is the simplicial cup product $(f\cup f')[v_0,v_1,v_2]=h_W\big(f[v_0,v_1],f'[v_1,v_2]\big)$ \cite[\S3.2]{Hatcher}. Equivariant cochains contribute the same term on every lift of a triangle, so dividing by the degree gives a sum over the triangles of $\PP^1$. Choose the lifts of the two triangles at $p_i$ so that they share their edge $g_i$; their edges $f_i$ are then two lifts that differ by the local monodromy at $p_i$, and the two triangles enter the fundamental class with opposite signs. The terms $h_W(f[b_0,p_i],f'[p_i,b_\infty])$ of the two triangles combine, by the cocycle condition on each triangle, into
\[
\mathcal S_{\rm simp}(f,f')=\sum_ih_W\big(v^f_{e_i}-v^f_{e_{i-1}},\ v^{f'}_{g_i}\big),
\]
with $v^f_\epsilon$ the value of $f$ on the chosen lift of $\epsilon$. The cocycle conditions on the two triangles at the $i$-th puncture give $a_i=Q_i^{-1}v_{f_i}$ and $v_{e_{i-1}}=v_{e_0}-F_i$. Substituting and telescoping gives $\mathcal S_{\rm simp}(f,f')=-\sum_i[h_W(c_i,a'_i)+h_W(Q_ic_i,F'_i)]$, the remaining term vanishing by the cocycle relation. By unitarity this equals $-\overline{\mathcal S(b,a)}$, and the identity $h_W(c_i,c'_i)=h_W(a_i,c'_i)+h_W(c_i,a'_i)$ together with telescoping $h_W(F_{i+1},F'_{i+1})-h_W(F_i,F'_i)$ gives $\mathcal S(a,b)+\overline{\mathcal S(b,a)}=0$. A coboundary pairs to zero by the same telescoping. Finally, replacing $a_i$ by $a_i+k$ with $T_ik=k$ changes neither $c_i$ nor any $F_j$, and changes $\mathcal S(a,b)$ by $h_W(k,c^b_i)=h_W(k,a^b_i)-h_W(T_i^{-1}k,a^b_i)=0$, by unitarity; the second variable enters only through the $c^b_i$.
\end{proof}

%% ------------------------------------------------------------------
\section{Dominance at a maximal degeneration}\label{sec:boundary}
%% ------------------------------------------------------------------

We give a criterion for dominance that only needs the monodromy tuple. The idea is to let the branch points collide in pairs until the cover degenerates as far as possible, to a point of the boundary of the Hurwitz space lying over a maximally degenerate stable rational curve. The local monodromies of the family around the boundary divisors there are unipotent after a base change, and their logarithms can be computed exactly from the tuple. If they span the full unipotent radical of the stabilizer of a rational Lagrangian, the nilpotent orbit theorem shows that the period map has full rank near the boundary. The same method was used by Alexeev, Donagi, Farkas, Izadi and Ortega for Prym--Tyurin maps to $\mathcal A_6$ \cite{ADFIO}.

\subsection{Twists and their logarithms}\label{subsec:twists}
Fix a tuple $\mathbf g$ in the component $S$ and $\Vp$, $\mathcal S$, $h$ as in Section~\ref{sec:herm}. Let $\gamma_1,\dots,\gamma_{r-3}$ be disjoint simple closed curves cutting the $r$-punctured sphere into pairs of pants. For each $j$:
\begin{enumerate}
\item[(i)] The full Dehn twist $\tau_j$ along $\gamma_j$ is a braid. If $\gamma_j$ encloses the consecutive points $a,\dots,b$ and $c=g_a\cdots g_b$, it acts on tuples by $g_i\mapsto cg_ic^{-1}$ for $a\le i\le b$, and equals $(\sigma_a\cdots\sigma_{b-1})^{b-a+1}$ in the Artin generators.
\item[(ii)] The \emph{return power} $k_j$ is the least $k\ge1$ with $\tau_j^k(\mathbf g)=\mathbf g$. The braid $\tau_j^{k_j}$ acts on $V_{\rm par}$ by a unitary operator $T_j$, computed by pulling back cocycles, as follows. Let $\rho\colon\pi_1(U_C,b_0)\to GL(V_\sigma)$ be $\sigma$ composed with the monodromy of the cover, so $\rho(\ell_i)=\sigma(g_i)$. A cocycle is a map $u$ with $u(xy)=u(x)+\rho(x)u(y)$; it is determined by $u_i=u(\ell_i)$, and it is parabolic when $u_i\in\operatorname{im}(1-\sigma(g_i))$ (so $u_i=c_i$ in the notation of Lemma~\ref{lem:cocycle}). Put $\gamma_c=\ell_a\cdots\ell_b$, so $\rho(\gamma_c)=\sigma(c)$, and
\[
u(\gamma_c)=\textstyle\sum_{i=a}^b\sigma(g_a\cdots g_{i-1})\,u_i .
\]
The twist maps $\ell_i$ to $\gamma_c\ell_i\gamma_c^{-1}$ for $a\le i\le b$ and fixes the other $\ell_i$, so the pulled-back cocycle $u'=u\circ\tau_j$ has
\[
u'_i=u(\gamma_c)+\sigma(c)\,u_i-\sigma(cg_ic^{-1})\,u(\gamma_c)\quad(a\le i\le b),\qquad u'_i=u_i\ \text{otherwise},
\]
by the cocycle identity and $u(\gamma^{-1})=-\rho(\gamma)^{-1}u(\gamma)$; it is a cocycle for the tuple $\tau_j(\mathbf g)$. The operator $T_j$ is the $k_j$-th iterate of this map, which preserves cocycles and coboundaries for $\mathbf g$.
\item[(iii)] Let $m_j$ be the order of the semisimple part of $T_j$; then $T_j^{m_j}$ is unipotent, and we put $N_j=\log T_j^{m_j}$.
\end{enumerate}
A twist with trivial boundary product $c$ fixes the tuple ($k_j=1$) but still acts on $\Vp$, in general nontrivially: in the example of Section~\ref{sec:g34} three such twists have logarithms of $K$-rank $4$.

Let $B$ be a connected complex manifold, $\phi\colon B\to\mathrm{Hur}$ a holomorphic map, and $\mathcal C\to B\times\PP^1$ the pulled-back family of $G$-covers, with $a\colon\mathcal C\to B$ and $\bar\pi\colon B\times\PP^1\to B$ the projections. Let $\mathcal U\subset B\times\PP^1$ be the complement of the branch divisor, $j\colon\mathcal U\to B\times\PP^1$ its inclusion, and $L'$ the local system on $\mathcal U$ defined fibrewise as in Section~\ref{sec:herm}. Fix $b\in B$ whose cover has monodromy tuple $\mathbf g$ with respect to a base point $b_0\in\PP^1$ and the loops of Lemma~\ref{lem:cocycle}, and assume that $b_0$ is not a branch point over any point of $B$. Fix a point $y_0$ of the cover over $(b,b_0)$. For a loop $\beta$ in $B$ at $b$, transporting $y_0$ along $\beta\times\{b_0\}$ in the $G$-torsor over $\mathcal U$ ends at $h_\beta y_0$ for a unique $h_\beta\in G$; and $\beta$ acts on the loops of $\PP^1-\{p_j\}$ by a braid, also written $\beta$.

\begin{lemma}\label{lem:gm}
The local system $\Hom_{K[G]}(V_\sigma,R^1a_*K)$ on $B$ is $R^1\bar\pi_*(j_*L')$, and its monodromy along $\beta$ is the pullback of parabolic cocycles by the braid $\beta$, as in (ii), composed with $\sigma(h_\beta)$; here $\beta(\mathbf g)=h_\beta\mathbf gh_\beta^{-1}$. If $\beta(\mathbf g)=\mathbf g$, then $h_\beta$ is central and $\sigma(h_\beta)$ is a root of unity in $K$; on the family of covers with a marked point over $b_0$, $h_\beta=1$.
\end{lemma}

\begin{proof}
Near a branch point a tame cover is $z\mapsto z^e$, so $p_*K=j_*p^\circ_*K$, and the finite map $p$ has no higher direct images; the Leray spectral sequence gives $R^1a_*K=R^1\bar\pi_*(p_*K)$, and taking the $\sigma$-multiplicity is exact. The description of the monodromy is \cite[Lemma~2.2 and Theorem~2.5]{DW1}, \cite[Theorem~1.3]{DW2}; on a family of $G$-covers the coefficient term comes from the $G$-torsor, so it is $\sigma(h_\beta)$. If $\beta(\mathbf g)=\mathbf g$, then $h_\beta$ centralizes a generating tuple, so it is central, and $\sigma(h_\beta)\in\End_{K[G]}(V_\sigma)=K$ by Schur's lemma.
\end{proof}

\begin{lemma}\label{lem:localmon}
Pinching the curves $\gamma_j$ gives a point of $\overline{\mathcal M}_{0,r}$ with $r-3$ nodes and smoothing parameters $q_1,\dots,q_{r-3}$. After the base change $w_j^{|Z(G)|\,k_jm_j}=q_j$, the family $A_\sigma$ has unipotent local monodromy around $w_j=0$, with logarithm $\varepsilon c_jN_j$ for positive integers $c_j$ and a sign $\varepsilon\in\{\pm1\}$ that does not depend on $j$.
\end{lemma}

\begin{proof}
The point lies on the normal-crossing boundary of $\overline{\mathcal M}_{0,r}$, with independent smoothing parameters \cite[Theorem~2.7]{Knudsen}. Over it the $G$-cover degenerates to an admissible cover: at a node the two inertia elements are inverse of the same order $e$, glued by the model $u=x^e$, $v=y^e$, $xy=t$, $uv=t^e$ \cite[\S4]{HarrisMumford}, after sufficiently divisible powers of the $q_j$. The loop of $q_j$ around $0$ is the full twist $\tau_j$ or its inverse. The local Hurwitz cover over the punctured polydisc is finite \cite[\S1.2]{FriedVolklein}, and the stabilizer of a sheet contains $\bigoplus_jk_j\Z$; so the Kummer cover in the statement dominates it, \'etale over the interior. By Lemma~\ref{lem:gm} the monodromy around $w_j=0$ is a power of $\zeta_jT_j$ with $\zeta_j^{|Z(G)|}=1$, raised to a multiple of $|Z(G)|\,m_j$, hence unipotent with logarithm a positive integer multiple of $\varepsilon_jN_j$, where $\varepsilon_j=\pm1$ according as the loop of $q_j$ is $\tau_j$ or $\tau_j^{-1}$.

It remains to see that $\varepsilon_j$ does not depend on $j$. In plumbing coordinates near the boundary point, the points enclosed by $\gamma_j$ lie, in a coordinate $z$ centred inside $\gamma_j$, at $z=q_jw$ with $w$ in a fixed compact set (for nested curves, at products of the $q_k$ of the curves containing them), and the points outside $\gamma_j$ do not depend on $q_j$. As $q_j$ runs once counterclockwise around $0$, with the other parameters fixed, the cluster inside $\gamma_j$ rotates rigidly once counterclockwise about its centre, for the complex orientation of $\PP^1$, and nothing outside $\gamma_j$ moves. This motion is the same kind of full twist for every $j$, so the loop of $q_j$ is $\tau_j^{\varepsilon}$ for one sign $\varepsilon$ independent of $j$, fixed by the orientation of Lemma~\ref{lem:cocycle} and the convention in (i).
\end{proof}

\subsection{The criterion}

Consider the following conditions on the logarithms $N_j$ of Lemma~\ref{lem:localmon}:
\begin{enumerate}
\item[(H1)] $N_jN_k=0$ for all $j,k$;
\item[(H2)] $I=\sum_j\operatorname{im}N_j$ is a $K$-subspace of dimension $n$, totally isotropic for $h$, with $N_j(I)=0$ for all $j$;
\item[(H3)] the Hermitian forms $\beta_j(x,y)=\mathcal S(x,N_jy)$, which descend to $\Vp/I$, are semidefinite of one common sign, and $\sum_j\beta_j$ is definite on $\Vp/I$;
\item[(H4)] the $N_j$ span a $\Q$-vector space of dimension $n^2$.
\end{enumerate}
By the several-variable nilpotent orbit theorem \cite[(4.12)]{Schmid}, \cite[(1.15)]{CKS}, after the base change of Lemma~\ref{lem:localmon} the lifted period map near the boundary point has the form $\tilde\Phi(z)=\exp(\sum_jz_jM_j)\cdot\psi(u)$ with $M_j=\varepsilon c_jN_j$, $u_j=e^{2\pi iz_j}$, and $\psi$ holomorphic on the polydisc with values in the compact dual. Put $F_0=\psi(0)$, the limit Hodge filtration. Fix an embedding $K\subset\C$ and write $\VpC$, $I_\C$ for the scalar extensions along it.

\begin{lemma}[The limit is Hodge--Tate]\label{lem:hodgetate}
Assume (H1)--(H3). For every $M=\sum_ja_jM_j$ with all $a_j>0$, $\operatorname{im}M=\ker M=I$, the weight filtration of $M$ is $W_0=W_1=I\subset W_2=\Vp$, the limit mixed Hodge structure is of Hodge--Tate type, and $\VpC=I_\C\oplus F^1_0$ with $\dim_\C F^1_0=n$.
\end{lemma}

\begin{proof}
By (H1) $M^2=0$, and by (H2) $\operatorname{im}M\subseteq I$ and $M(I)=0$. By (H3), for $x\notin I$ we have $\mathcal S(x,\varepsilon Mx)=\sum_ja_jc_j\beta_j(x,x)\ne0$, so $M\colon \Vp/I\to I$ is injective, hence bijective since both sides have $K$-dimension $n$. Thus $\operatorname{im}M=\ker M=I$, and the monodromy weight filtration of $M$, centred at $1$, is $W_0=W_1=I$, $W_2=\Vp$. By \cite[(4.5)]{CKS}, $(W,F_0)$ is a mixed Hodge structure polarized by $M$. Since the variation has weight one, $\operatorname{Gr}^W_0=I$ has type $(0,0)$ and $\operatorname{Gr}^W_2=\Vp/I$ type $(1,1)$. A Hodge structure of type $(0,0)$ has $F^1=0$, so $F^1_0\cap I_\C=0$. In the chosen eigenspace the weight-one variation has Hodge numbers $(n,n)$, so $\dim F^1_0=n=\dim I_\C$ and $\VpC=I_\C\oplus F^1_0$.
\end{proof}

\begin{lemma}[The logarithms span the tangent space]\label{lem:tangentspan}
Assume (H1)--(H4). Then $v\mapsto\sum_jv_jM_j\cdot F_0$ maps $\C^{\{j\}}$ onto the tangent space $\Hom(F^1_0,\VpC/F^1_0)$ of $\cHn$ at $F_0$, of dimension $n^2$.
\end{lemma}

\begin{proof}
The rational Lie algebra $\mathfrak n_I=\{X\in\mathfrak u(h):X(\Vp)\subseteq I,\ X(I)=0\}$, a $\Q$-vector space ($\mathfrak u(h)$ is the Lie algebra of the unitary group of $h$, an algebraic group over $\Q$), has dimension $n^2$ over $\Q$: via $X\mapsto\mathcal S(\cdot,X\cdot)$ it is the $\Q$-space of $K$-Hermitian forms on the $n$-dimensional $K$-space $\Vp/I$, which has $\Q$-dimension $n+2\binom n2=n^2$ ($n$ rational diagonal entries and $\binom n2$ entries in $K$ above the diagonal). Its complexification maps isomorphically onto $\Hom(\VpC/I_\C,I_\C)$. By (H1)--(H2) every $N_j$ lies in $\mathfrak n_I$, and by (H4) they span it. By Lemma~\ref{lem:hodgetate}, $\VpC/I_\C\cong F^1_0$ and $I_\C\cong \VpC/F^1_0$, so an element $X$ of $\mathfrak n_I\otimes\C$ is determined by $X|_{F^1_0}$, and $X\cdot F_0$, the class of $X|_{F^1_0}$ in $\Hom(F^1_0,\VpC/F^1_0)$, vanishes only if $X(F^1_0)\subseteq F^1_0\cap I_\C=0$. So $X\mapsto X\cdot F_0$ is injective on the $n^2$-dimensional space $\mathfrak n_I\otimes\C$, hence onto the $n^2$-dimensional tangent space.
\end{proof}

\begin{proposition}\label{prop:boundary}
Assume (H1)--(H4). Then $S$ is dominant, and the discriminant $\fd(S)$ is split.
\end{proposition}

\begin{proof}
\emph{The period map near the boundary.} After translating by $\exp(-\sum_jz_jM_j)$, the $j$-th column of the derivative of $\tilde\Phi$ is $M_j\cdot\psi(u)+2\pi iu_j\,\partial\psi/\partial u_j$, which tends to $M_j\cdot F_0$ as $u\to0$. By Lemma~\ref{lem:tangentspan} these limits span the $n^2$-dimensional tangent space, and so we may choose $n^2$ of the $r-3$ indices $j$ for which the limits $M_j\cdot F_0$ form a basis; the corresponding $n^2\times n^2$ minor of the limit is nonzero and stays nonzero for small $u$, so the period map has rank $n^2$ at interior points near the boundary, and $S$ is dominant.

\emph{The discriminant.} By (H2), $(\Vp,h)$ has a totally isotropic subspace of dimension $n$, so $h$ is split, and so is the Hermitian form of $A_\sigma(C)$ by Proposition~\ref{prop:herm}.
\end{proof}

\begin{remark}\label{rem:boundary}
By (H2) and Proposition~\ref{prop:boundary}, the criterion applies only to split discriminants; for context, Brosnan \cite{Brosnan} shows more generally that a family of abelian $2n$-folds of Weil type with a maximal degeneration, one whose monodromy logarithm has rank $2n$ over $\Q$, has split Hermitian form. In the example of Section~\ref{sec:g34} the parabolic cohomology of every pair of pants vanishes, as it must: by the Clemens--Schmid sequence \cite{Clemens} the weight-one part of the limit is the sum, over the components of the degenerate curve, of their parabolic cohomology with coefficients in $L'$, and the limit is of Hodge--Tate type (Lemma~\ref{lem:hodgetate}).
\end{remark}

%% ------------------------------------------------------------------
\section{The group $G_{34}$ and the proof of Theorem~A}\label{sec:g34}
%% ------------------------------------------------------------------

Let $K=\Q(\sqrt{-3})$ and $\omega$ a primitive cube root of unity. The Shephard--Todd group $G_{34}\subset GL_6(K)$ \cite{ShephardTodd} has order $39\,191\,040$ and centre $\mu_6$. Let $\sigma$ be its reflection representation on $K^6$, of dimension $d=6$ and character field $K$; it preserves the standard Hermitian form $\langle x,y\rangle=\sum_k\bar x_ky_k$. Its reflections have order two and form one conjugacy class.

Let $\mathbf g=(s_1,s_1,s_2,s_2,\dots,s_{10},s_{10})$ be the tuple of twenty reflections listed in Appendix~\ref{app:data}, with $s_i$ the reflection in the root $v_i$. Its product is $1$ because every $s_i$ is an involution, and the $s_i$ generate $G_{34}$ (Appendix~\ref{app:data}). So $\mathbf g$ lies in the Nielsen class of twenty reflections; let $S$ be its component of the Hurwitz space, of dimension $17$.

\emph{Why this tuple and these curves.} With $r$ reflections the count of Lemma~\ref{lem:g34mult} gives $m_\sigma=r-12=2n$, so eightfolds ($n=4$) need exactly $r=20$, and then the Hurwitz space has dimension $r-3=17\ge n^2=16$, which leaves room for dominance. One step further fails: $r=22$ gives $n=5$ with $19<25$ moduli. We need eightfolds rather than sixfolds because the boundary criterion only reaches split discriminants, and Section~\ref{sec:compensation} turns split eightfolds into all sixfolds. Doubling each reflection makes the product of the tuple trivial with no computation; that the ten distinct reflections still generate $G_{34}$ is checked in Appendix~\ref{app:data}. Doubling also makes many boundary curves easy: a curve enclosing a pair $s_i,s_i$, or a block of complete pairs, has trivial boundary product, so its twist fixes the tuple ($k_j=1$) and its logarithm comes from a single pullback. The seventeen curves of Appendix~\ref{app:data} are nested or disjoint, so they cut the twenty-punctured sphere into pairs of pants, a maximal degeneration ($r-3=17$ curves). They were found by a search for a decomposition satisfying the conditions of Proposition~\ref{prop:boundary}; other components of the Nielsen class and other pants decompositions were not examined.

\begin{lemma}\label{lem:g34mult}
The representation $\sigma$ occurs in $H^1(C)$ with multiplicity $m_\sigma=8$, with Hodge numbers $(4,4)$. So $A_\sigma(C)$ is an abelian eightfold of Weil type for $K$.
\end{lemma}

\begin{proof}
A reflection has $c_\sigma=1$, and its eigenvalues are $-1$ once and $1$ five times, so $\age\sigma(g_j^{\pm1})=\tfrac12$. Proposition~\ref{prop:mult} gives $m_\sigma=-12+20=8$ and $20\cdot\tfrac12-6=4$ for each Hodge number; conclude with Lemma~\ref{lem:weiltype}.
\end{proof}

\begin{proof}[Proof of Theorem~A]
The reflections of $G_{34}$ are involutions, so the branch classes are real, and condition~(L) holds by Theorem~\ref{thm:brauer}. Take the $17$ curves of Appendix~\ref{app:data}, each enclosing consecutive points of the twenty. For them the operators $N_j$ of Section~\ref{sec:boundary}, computed exactly over $K$ from $\mathbf g$ by Lemma~\ref{lem:cocycle} and the pullback formula of Section~\ref{sec:boundary}, satisfy (H1)--(H4): $N_jN_k=0$; the images span a four-dimensional $K$-subspace $I$, totally isotropic for $h$ and killed by every $N_j$; every $\beta_j$ is negative semidefinite and $\sum_j\beta_j$ is negative definite on $\Vp/I$; and the $N_j$ span a $\Q$-space of dimension $16=n^2$. Appendix~\ref{app:data} lists the data. By Proposition~\ref{prop:boundary}, $S$ is dominant and $\fd(S)$ is split. Theorem~\ref{thm:transfer} gives the Weil classes of every polarized abelian eightfold of Weil type for $K$ with split discriminant.
\end{proof}

\begin{remark}\label{rem:witness}
Theorem~\ref{thm:brauer} does not produce the coefficients, but an explicit identity is also available. Here the $120$ non-real irreducible representations $\tau$ of $G_{34}$ with $m_\tau>0$ form $60$ pairs, and no quaternionic $\tau$ has $m_\tau>0$, so $X^*(T)\cong\Z^{60}$; and $e_\sigma=\sum_in_ia(\chi_i)$ with $84$ admissible suppliers, on cyclic subgroups and on subgroups $\langle x\rangle\times\operatorname{Syl}_p(C_G(x))$, $p=2,3$, with $\sum_i|n_i|=231$. The supplementary files contain the identity and a script that rebuilds every row from its subgroup and character, with its $q$, $s$, $h$ and admissibility. Three of the suppliers have $m=12$ and $s>0$, so the identity uses Lemma~\ref{lem:schoenR}. Cyclic subgroups alone generate a sublattice of full rank that does not contain $e_\sigma$: the smallest $k$ with $ke_\sigma$ in it is $31\,104=2^7\cdot3^5$, and explicit linear functionals modulo $2$ (on six coordinates) and modulo $3$ (on four) vanish on the cyclic sublattice and take the value $1$ on $e_\sigma$.
\end{remark}

\begin{example}[One supplier of the identity]\label{ex:supplier}
The term of the identity in Remark~\ref{rem:witness} with the smallest $h$ is the following. Let $H'=\langle-I_6\rangle\times P$, where $P$ is a Sylow $3$-subgroup of $G_{34}$, so $|H'|=2\cdot3^7=4\,374$; $H'$ is elementary in Brauer's sense. Let $\chi$ be a linear character of $H'$ of order $m=3$, trivial on $-I_6$ (the character of row $54$ of the deposited rebuild log). The genus of $C$ follows from Riemann--Hurwitz for $C\to\PP^1$: each of the twenty branch points has $|G_{34}|/2$ preimages, so $2g(C)-2=|G_{34}|(-2+20\cdot\tfrac12)=8|G_{34}|$. A point of $C$ fixed by a nontrivial element of $G_{34}$ lies over a branch point, and its stabilizer is generated by a reflection. The only involution in $H'$ is $-I_6$, since $|P|$ is odd, and $-I_6$ is not a reflection; so $H'$ acts freely on $C$, and both $C\to C/H'$ and $C/\ker\chi\to C/H'$ are unramified: $s=0$. Hence $2g(C)-2=|H'|(2q-2)$, that is $8\cdot39\,191\,040=4\,374\,(2q-2)$, so $q=35\,841$ and $h=2q-2=71\,680$, even and positive; with $s=0$ the supplier is admissible (Lemma~\ref{lem:schoen}). It occurs in the identity with coefficient $3$, and it is itself a Brauer term, since $H'$ is elementary and $\chi$ linear.
\end{example}

%% ------------------------------------------------------------------
\section{From eightfolds to sixfolds: proof of Theorem~B}\label{sec:compensation}
%% ------------------------------------------------------------------

The argument is Schoen's \cite[Proposition~10]{Schoen98}, in the form of Voisin \cite[Lemme~2.9]{Voisin}: if a product $A\times Y$ has algebraic Weil classes and $Y$ has algebraic Weil classes, so does $A$. We need only a factor $Y$ that changes the discriminant to the split one. It works for every imaginary quadratic $K$. Let $E$ be an elliptic curve with complex multiplication by $\mathcal O_K$, and $\bar E$ the same curve with the conjugate action of $K$.

\begin{lemma}\label{lem:eebar}
For an integer $m\ge1$, the surface $E\times\bar E$ with the product polarization of degrees $(1,m)$ is of Weil type for $K$, of signature $(1,1)$ and discriminant $[-m]$, and its Weil classes are divisor classes.
\end{lemma}

\begin{proof}
The two factors carry the two embeddings of $K$ on their tangent spaces: $\iota(\vartheta)$ acts on $T_0E$ by $\vartheta$ and on $T_0\bar E$ by $\bar\vartheta=-\vartheta$. A polarization $E_0$ of a curve satisfies $E_0(x,Jx)>0$ for its complex structure $J$, so the values $H(x,x)=E_0(x,\iota(\vartheta)x)$ have opposite signs on the two factors. Identifying $H_1(\bar E,\Q)$ with $H_1(E,\Q)$, the polarization of degree $m$ is $m$ times that of degree $1$, so the two rank-one Hermitian forms are $(c)$ and $(-cm)$ with the same $c\in\Q^\times$, and the sum has determinant $-c^2m$, of class $[-m]$. The Weil classes lie in $H^2$ and have type $(1,1)$, so they are divisor classes by the Lefschetz $(1,1)$ theorem.
\end{proof}

\begin{proposition}\label{prop:compensation}
Let $n\ge4$, and suppose the Weil classes of every polarized abelian $2n$-fold of Weil type for $K$ with split discriminant are algebraic. Then the Weil classes of every polarized abelian sixfold of Weil type for $K$ are algebraic.
\end{proposition}

\begin{proof}
\emph{The product.} Let $A$ be a sixfold of discriminant $[t]$, $t<0$. Put $k=n-3$ and $Y=\prod_{j=1}^kY_j$ with $Y_j=E\times\bar E$ polarized with degrees $(1,m_j)$, where $m_2=\dots=m_k=1$ and $m_1$ is a positive integer in the class $[(-1)^{k+n}/t]$; that class is positive, and a positive class has a positive integral representative because squares are norms. By Lemma~\ref{lem:eebar}, $A\times Y$ is of Weil type for $K$, of signature $(n,n)$ and discriminant $[t]\prod_j[-m_j]=[t]\,[-m_1]\,[-1]^{k-1}=[(-1)^n]$, the split class; so its Weil classes are algebraic. In the case used for Theorem~B, $n=4$ and $k=1$: $Y=E\times\bar E$ with degrees $(1,m)$, $m\in[-1/t]$, and $[t][-m]=[1]$.

\emph{Notation.} Fix an embedding $\iota$ of $K$ in $\C$ and write $\bar\iota$ for the other. For a $K$-module $M$ of Weil type, $W_K(M)\otimes\C$ is the sum of two lines, on which $a\in K$ (acting through $\iota$ on $H^1$) acts on the top exterior power by $\iota(a)^{\rm top}$ and $\bar\iota(a)^{\rm top}$; we call them the $\iota$-line and the $\bar\iota$-line, and complex conjugation exchanges them. Let $x$ span the $\iota$-line of $W_K(A)\otimes\C$, so $\bar x$ spans the $\bar\iota$-line, and let $w_j$, $\bar w_j$ span those of $W_K(Y_j)\otimes\C\subset H^2(Y_j,\C)$. Since $\bigwedge^{2n}_KH^1(A\times Y)=\bigwedge^6_KH^1(A)\otimes\bigotimes_j\bigwedge^2_KH^1(Y_j)$, a nonzero rational Weil class of $A\times Y$ has the form $\xi=b\,x\otimes\bigotimes_jw_j+\bar b\,\bar x\otimes\bigotimes_j\bar w_j$ with $b\in\C^\times$.

\emph{The pairings on $Y_j$.} An element $a\in\mathcal O_K$ acts on $Y_j$ by an endomorphism of degree $\Nm(a)^2$, on the $\iota$-line of $H^2$ by $\iota(a)^2$ and on $H^4(Y_j)$ by $\Nm(a)^2=\iota(a)^2\bar\iota(a)^2$. So $w_j\cup w_j$, on which $a$ acts by $\iota(a)^4$, vanishes (take $a$ with $\iota(a)^2\notin\R$), and by Poincar\'e duality $\epsilon_j=\int_{Y_j}w_j\cup\bar w_j\ne0$.

\emph{The contraction.} A nonzero rational class $y_j=a_jw_j+\bar a_j\bar w_j\in W_K(Y_j)$ is algebraic (Lemma~\ref{lem:eebar}). Then
\[
z=\operatorname{pr}_{A*}\big(\xi\cup\operatorname{pr}_Y^*(y_1\boxtimes\dots\boxtimes y_k)\big)=\pm\Big(b\textstyle\prod_j\bar a_j\epsilon_j\Big)\,x+\pm\overline{\Big(b\textstyle\prod_j\bar a_j\epsilon_j\Big)}\,\bar x,
\]
with one common sign, a nonzero rational algebraic class in $W_K(A)$.

\emph{Two independent classes.} For $\alpha\in\mathcal O_K$, the class $\alpha^*y_1$ is algebraic, and it multiplies $\bar a_1$ by $\bar\iota(\alpha)^2$; so the class $z_\alpha$ obtained from $\alpha^*y_1$ has coordinates $(c\,\bar\iota(\alpha)^2,\ \overline{c\,\bar\iota(\alpha)^2})$ in the basis $x,\bar x$, with $c\ne0$ independent of $\alpha$. The determinant of the coordinates of $z_1$ and $z_\alpha$ is
\[
c\,\overline{c\,\bar\iota(\alpha)^2}-\bar c\,c\,\bar\iota(\alpha)^2=|c|^2\big(\iota(\alpha)^2-\bar\iota(\alpha)^2\big),
\]
which is nonzero exactly when $\alpha^2\notin\Q$. Taking $\alpha\in\mathcal O_K$ neither rational nor purely imaginary (for instance $\alpha=1+\omega$ when $K=\Q(\sqrt{-3})$, with $\alpha^2=\omega$), $z_1$ and $z_\alpha$ are two independent algebraic classes in the two-dimensional $W_K(A)$.
\end{proof}

\begin{proof}[Proof of Theorem~B]
Proposition~\ref{prop:compensation} with $n=4$ and Theorem~A.
\end{proof}

%% ------------------------------------------------------------------
\section{Powers: proof of Theorem~C}\label{sec:powers}
%% ------------------------------------------------------------------

Let $W$ be a complex vector space of dimension $N$, and $U,U'$ complex vector spaces of dimension $m$. Put $\Lambda=\bigwedge(W\otimes U)\otimes\bigwedge(W^*\otimes U')$, with the commuting actions of $GL(W)$ and $GL(U)\times GL(U')$. For $u\in U$, $u'\in U'$ let $\theta(u,u')=\sum_i(w_i\otimes u)\wedge(w_i^*\otimes u')$ for dual bases, and for $a\in U$ let $\det(W\otimes a)\in\bigwedge^N(W\otimes U)$ be the image of a generator of $\bigwedge^NW$ under $w\mapsto w\otimes a$; define $\det(W^*\otimes a')$ similarly.

\begin{lemma}\label{lem:sl}
$\Lambda^{SL(W)}$ is generated by the $\theta(u,u')$, the $\det(W\otimes a)$ and the $\det(W^*\otimes a')$.
\end{lemma}

\begin{proof}
We deduce it from the first fundamental theorem of invariant theory for $SL(W)$, in the form of Kraft and Procesi \cite[\S8.4, Theorem (FFT for $SL_n$), p.~90]{KraftProcesi}: \emph{the ring of invariants $K[V^p\oplus V^{*q}]^{SL(V)}$ is generated by the scalar products $\langle j\mid i\rangle$ and the determinants $[i_1,\dots,i_n]$ and $[j_1,\dots,j_n]^*$}, for a vector space $V$ of dimension $n$ over a field $K$ of characteristic zero (see also \cite[Ch.~II]{Weyl}). We use it in multilinear form. An $SL(W)$-invariant element of $W^{\otimes p}\otimes W^{*\otimes q}$ is an invariant multilinear function on $W^{*p}\times W^q$, that is, an invariant polynomial on $W^{*p}\oplus W^q$ that is homogeneous of degree one in each of the $p+q$ vector variables. Apply the theorem with $V=W^*$. A polynomial in the scalar products and determinants that is homogeneous of degree one in each variable is a linear combination of products of scalar products and determinants in which every variable occurs exactly once. So the invariants in $W^{\otimes p}\otimes W^{*\otimes q}$ are spanned by tensor products of three kinds of elementary invariants, each occupying some of the tensor factors: the contraction $\sum_iw_i\otimes w_i^*$ in one factor $W$ and one factor $W^*$; the determinant tensor $\sum_{\pi\in S_N}\operatorname{sgn}(\pi)\,w_{\pi(1)}\otimes\dots\otimes w_{\pi(N)}$ in $N$ factors $W$; and the dual determinant tensor in $N$ factors $W^*$. (This is the skew $(GL,GL)$ duality of \cite{Howe} in another form.)

The homogeneous piece $\bigwedge^p(W\otimes U)\otimes\bigwedge^q(W^*\otimes U')$ of $\Lambda$ is the image of $(W\otimes U)^{\otimes p}\otimes(W^*\otimes U')^{\otimes q}=W^{\otimes p}\otimes W^{*\otimes q}\otimes U^{\otimes p}\otimes U'^{\otimes q}$ under antisymmetrization, which is $SL(W)$-equivariant and surjective. Since $SL(W)$ is reductive, the invariants of the image are the image of the invariants, so they are spanned by the images of $I\otimes u_1\otimes\dots\otimes u_p\otimes u'_1\otimes\dots\otimes u'_q$ with $I$ a product of elementary invariants. Antisymmetrization turns tensor products into wedge products, up to sign. The image of a contraction in the factors $(w\otimes u_k)$ and $(w^*\otimes u'_l)$ is $\theta(u_k,u'_l)$. The image of a determinant tensor in the factors $u_{k_1},\dots,u_{k_N}$ is $\sum_\pi\operatorname{sgn}(\pi)(w_{\pi(1)}\otimes u_{k_1})\wedge\dots\wedge(w_{\pi(N)}\otimes u_{k_N})$, which is the full polarization of the homogeneous polynomial map $a\mapsto N!\det(W\otimes a)$ of degree $N$, evaluated at $(u_{k_1},\dots,u_{k_N})$; by the polarization identity it is a linear combination of values $\det(W\otimes a)$. The same holds for the dual determinant tensors. So every invariant is a combination of products of the stated generators.
\end{proof}

\begin{theorem}\label{thm:powers}
Let $(A,\iota,E)$ be a polarized abelian $2n$-fold of Weil type for $K$, $n\ge2$, whose Weil classes are algebraic and whose Hodge group is $SU_H$. Then every Hodge class on every power $A^m$ is algebraic.
\end{theorem}

\begin{proof}
Let $H^1(A,\C)=W\oplus W'$ be the decomposition into the two eigenspaces of $K$. The polarization identifies $W'$ with $W^*$ as representations of $SU_H(\C)=SL(W)$. The Hodge group of $A^m$ is $SU_H$ acting diagonally, so the Hodge classes on $A^m$, tensored with $\C$, are $\Lambda^{SL(W)}$ with $N=2n$ and $U=U'=\C^m$. The $\theta(u,u')$ are Hodge classes of degree two, hence combinations of divisor classes. Let $\mathcal O=\{a\in K:\iota(a)\in\End(A)\}$. For $a\in\mathcal O^m$ the homomorphism $A^m\to A$, $(x_i)\mapsto\sum_i\iota(a_i)x_i$, pulls the Weil classes of $A$ back to $\det(W\otimes\tilde a)$ and $\det(W^*\otimes\overline{\tilde a})$, with $\tilde a$ the image of $a$ under the given embedding; these are algebraic. Since $\mathcal O$ is a lattice in $K$, the set of $\tilde a$ is Zariski dense in $U$, so these classes span all $\det(W\otimes a)$, and likewise for $W^*$. By Lemma~\ref{lem:sl}, products of algebraic classes span all Hodge classes over $\C$, and rational descent as in Corollary~\ref{cor:everymember} finishes the proof.
\end{proof}

\begin{proof}[Proof of Theorem~C]
Theorem~B and Theorem~\ref{thm:powers} with $n=3$. For sixfolds the hypothesis $\Hg(A)=SU_H$ holds outside a countable union of proper closed subvarieties of the moduli space \cite[Theorem~6.11]{vanGeemen}.
\end{proof}

%% ------------------------------------------------------------------
\section{Abelian sixfolds of type III}\label{sec:typeIII}
%% ------------------------------------------------------------------

Theorem~B covers every member of every moduli space, including special members, and this gives a result for abelian sixfolds of type III.

\begin{corollary}\label{cor:typeIII}
Let $D$ be a totally definite quaternion algebra over $\Q$ containing $\Q(\sqrt{-3})$, that is, ramified at $\infty$ and at an odd number of primes, each equal to $3$ or congruent to $2$ modulo $3$. Let $\mathcal A$ be a family of abelian sixfolds of PEL type for $D$, that is, the universal family over a connected Shimura variety parametrizing polarized abelian sixfolds with an action of a fixed order of $D$, whose general member has endomorphism algebra exactly $D$ \cite[\S2.5]{Abdulali}. For every member $A$ of $\mathcal A$ whose Hodge group is the generic Hodge group of the family, every Hodge class on every power of $A$ is algebraic.
\end{corollary}

\begin{proof}
$K=\Q(\sqrt{-3})$ embeds in $D$ as a maximal subfield \cite[Proposition~14.6.7]{Voight}. Restricting the action of $D$ to $K$ makes $A$ an abelian sixfold of Weil type for $K$ with a compatible polarization \cite[\S4]{Abdulali}: an element of $D$ anticommuting with $\sqrt{-3}$ exchanges the two eigenspaces of $K$ on $H^{1,0}(A)$. By Theorem~B, $W_K(A)$ is spanned by algebraic classes.

Abdulali's Theorem~4.1 \cite[p.~673]{Abdulali} states: \emph{Let $A$ be a general member of an effective PEL-family of type III. The Hodge ring of $A$ is generated by divisors, together with the Weil cycles associated to a maximal subfield of $\End_\Q(A)$. The algebraicity of the Weil cycles on $A$ implies the usual Hodge conjecture for all powers of $A$.} In his \S4, ``general'' is used for a fibre whose Hodge group is the group $G_1$ defined there, the generic Hodge group of the family: every fibre has Hodge group contained in $G_1$, with equality for the general fibre. His proof computes the Hodge ring of $A^r$ as the ring of invariants of $G_1$ and uses nothing else about $A$, so it applies verbatim to every member whose Hodge group is $G_1$. For such a member, $\End^0(A)$ is the commutant of $G_1$, which is $D$, so $K$ is a maximal subfield of $\End^0(A)$ and the Weil cycles of the theorem may be taken to be those for $K$, which are algebraic.
\end{proof}

\begin{example}
The point of Corollary~\ref{cor:typeIII} is that type III sixfolds reach non-split discriminants, so it is not a consequence of the split case. Write $D=(a,b)_\Q$ with basis $1,u,v,uv$, $u^2=a$, $v^2=b$, and identify $K=\Q(u)$. A polarization of a sixfold of type III is the trace of a skew-Hermitian form $T(x,y)=\bar x^{T}\,\mathrm{diag}(t_1,t_2,t_3)\,y$ on $D^3$ with pure quaternions $t_k$ \cite[\S2.5]{Abdulali}. As a right $K$-space, $D=K\oplus vK$, and the $K$-component of $T$ is the skew-Hermitian form of Section~\ref{sec:herm} whose trace is the polarization. For one coordinate with $t=\alpha u+\beta v+\gamma uv$, its Gram matrix in the basis $1,v$ is
\[
\begin{pmatrix}\alpha u& b(\beta+\gamma u)\\-b(\beta-\gamma u)&\alpha bu\end{pmatrix},
\]
of determinant $b(a\alpha^2+b\beta^2-ab\gamma^2)=-b\,\operatorname{nrd}(t)$. Dividing by $u$ to make it Hermitian multiplies the determinant by $1/a$. So the discriminant of $A$ is the class of $\prod_k(-b\,\operatorname{nrd}(t_k)/a)$ in $\Q^\times/\Nm K^\times$. For $a=-3$ this is $[b^3\prod_k\operatorname{nrd}(t_k)]$, since $3=\Nm(\sqrt{-3})$. Take $D=(-3,-1)_\Q$, so $\operatorname{nrd}(\alpha u+\beta v+\gamma uv)=3\alpha^2+\beta^2+3\gamma^2$. The forms $\mathrm{diag}(v,u,u+2uv)$ and $\mathrm{diag}(u+4v+uv,u,u+uv)$ have $\prod_k\operatorname{nrd}(t_k)=1\cdot3\cdot15=45$ and $22\cdot3\cdot6=396$, so discriminants $[-45]=[-5]$ and $[-396]=[-11]$. The split class is $[(-1)^3]=[-1]$, and $[-5]=[-1]$ would make $5$ a norm from $K$, which fails in $\Q_5$: since $5$ is inert in $K$, the completion $K\otimes\Q_5$ is the unramified quadratic extension of $\Q_5$, whose norms have even $5$-adic valuation, and $v_5(5)=1$ is odd. The same holds for $11$. So both discriminants are non-split.
\end{example}

\appendix
%% ------------------------------------------------------------------
\section{The data for $G_{34}$}\label{app:data}
%% ------------------------------------------------------------------

\emph{The tuple.} Write elements of $K=\Q(\omega)$ as $a+b\omega$. The reflections $s_i(x)=x-2\frac{\langle v_i,x\rangle}{\langle v_i,v_i\rangle}v_i$ in the roots
\[
\setlength{\arraycolsep}{2pt}\small
\begin{array}{ll}
v_1=(1,\,-1-\omega,\,1,\,\omega,\,-1-\omega,\,\omega), & v_6=(1,\,-1-\omega,\,\omega,\,\omega,\,\omega,\,\omega),\\
v_2=(1,\,0,\,0,\,1+\omega,\,0,\,0), & v_7=(0,\,0,\,1,\,0,\,0,\,-1),\\
v_3=(0,\,1,\,1+\omega,\,0,\,0,\,0), & v_8=(1,\,\omega,\,\omega,\,-1-\omega,\,1,\,-1-\omega),\\
v_4=(1,\,\omega,\,-1-\omega,\,\omega,\,-1-\omega,\,1), & v_9=(1,\,-\omega,\,0,\,0,\,0,\,0),\\
v_5=(1,\,-1-\omega,\,-1-\omega,\,-1-\omega,\,-1-\omega,\,\omega), & v_{10}=(0,\,1,\,0,\,0,\,0,\,-1),
\end{array}
\]
with $\langle\ ,\ \rangle$ the standard Hermitian form, generate a group of order $39\,191\,040$ with centre of order $6$. In GAP, with \texttt{w:=E(3)}, the ten matrices $1-2\,v_i\bar v_i^{\,T}/\langle v_i,v_i\rangle$ give these numbers through \texttt{Size(Group(gens))} and \texttt{Size(Centre(Group(gens)))}. By the classification of Shephard and Todd \cite{ShephardTodd}, the only irreducible complex reflection group in dimension $6$ of this order is $G_{34}$: the imprimitive groups $G(m,p,6)$, $m\ge2$, have order $720\,m^6/p$, which never equals $39\,191\,040$; the symmetric group $S_7=G(1,1,7)$, acting on its six-dimensional reflection representation, has order $5040$; and the only other primitive group in dimension $6$, the Weyl group $G_{35}$ of $E_6$, has order $51\,840$; the group acts irreducibly on $K^6$, because its reduction modulo the prime $\mathfrak p=(7,\omega-2)$ of $\mathcal O_K$ is absolutely irreducible over $\mathcal O_K/\mathfrak p=\mathbb F_7$ (checked by \texttt{MTX.IsAbsolutelyIrreducible} on the reduced matrices). Here is why that suffices. Every $\langle v_i,v_i\rangle$ is $2$ or $6$, so the matrices $1-2v_i\bar v_i^{\,T}/\langle v_i,v_i\rangle$ have entries in $\mathcal O_K[\tfrac13]$, which is contained in the local ring $\mathcal O_{\mathfrak p}$; so $\Lambda=\mathcal O_{\mathfrak p}^6$ is a stable lattice and the reduction is $\Lambda/\mathfrak p\Lambda$. If $0\ne W\subsetneq K^6$ were a stable subspace, then $W\cap\Lambda$ would be a saturated stable sublattice, and its image in $\Lambda/\mathfrak p\Lambda$ would be a stable subspace of dimension $\dim_KW$, neither $0$ nor everything. The tuple is $\mathbf g=(s_1,s_1,s_2,s_2,\dots,s_{10},s_{10})$.

\emph{The curves.} The seventeen curves $\gamma_j$ each enclose the consecutive points $a,\dots,b$ of the twenty. The table gives, for each, the return power $k_j$, the order $m_j$ of the semisimple part of $T_j$, the $K$-rank of $N_j$, and the signature of $\beta_j$ on $\Vp/I$ as (positive, negative, zero).
\begin{center}
\begin{tabular}{cccccc}
$j$ & $[a,b]$ & $k_j$ & $m_j$ & $\rk N_j$ & signature of $\beta_j$\\\hline
1 & $[1,8]$ & 1 & 1 & 4 & $(0,4,0)$\\
2 & $[1,18]$ & 1 & 1 & 1 & $(0,1,3)$\\
3 & $[2,8]$ & 2 & 1 & 3 & $(0,3,1)$\\
4 & $[3,8]$ & 1 & 1 & 3 & $(0,3,1)$\\
5 & $[4,8]$ & 2 & 1 & 2 & $(0,2,2)$\\
6 & $[5,6]$ & 1 & 1 & 1 & $(0,1,3)$\\
7 & $[5,8]$ & 1 & 1 & 2 & $(0,2,2)$\\
8 & $[7,8]$ & 1 & 1 & 1 & $(0,1,3)$\\
9 & $[9,10]$ & 1 & 1 & 1 & $(0,1,3)$\\
10 & $[9,12]$ & 1 & 1 & 2 & $(0,2,2)$\\
11 & $[9,16]$ & 1 & 1 & 4 & $(0,4,0)$\\
12 & $[9,17]$ & 2 & 1 & 4 & $(0,4,0)$\\
13 & $[9,18]$ & 1 & 1 & 4 & $(0,4,0)$\\
14 & $[11,12]$ & 1 & 1 & 1 & $(0,1,3)$\\
15 & $[13,14]$ & 1 & 1 & 1 & $(0,1,3)$\\
16 & $[13,16]$ & 1 & 1 & 2 & $(0,2,2)$\\
17 & $[15,16]$ & 1 & 1 & 1 & $(0,1,3)$
\end{tabular}
\end{center}
Fourteen of the curves have trivial boundary product in $G_{34}$; for the three with $k_j=2$ the tuple first returns after two twists. The sum $\sum_j\beta_j$ has signature $(0,4,0)$. The $N_j$ satisfy $N_jN_k=0$, their images span a four-dimensional totally isotropic $K$-subspace $I$ killed by every $N_j$, and they span a $\Q$-space of dimension $16$. Every pair of pants has vanishing parabolic cohomology (Remark~\ref{rem:boundary}).

\emph{The basis.} A parabolic cocycle $u$ for $\mathbf g$ is determined by its values $u_i=u(\ell_i)$, $i=1,\dots,20$, and $u_i$ lies in the image of $1-s_{\lceil i/2\rceil}$, which is the line $Kv_{\lceil i/2\rceil}$; so $u_i=\lambda_iv_{\lceil i/2\rceil}$ for scalars $\lambda_i\in K$. We use the basis $b_1,\dots,b_8$ of $\Vp$ given by the following cocycles; all $\lambda_i$ not listed are $0$. For $k=1,\dots,4$, $b_k$ has $\lambda_{2k-1}=\lambda_{2k}=1$, that is, the value $v_k$ on both loops around the two points with monodromy $s_k$. For $k=5,\dots,8$:
\begin{center}
\small
\begin{tabular}{c|cccccccccc}
$i$ & $1$ & $3$ & $5$ & $7$ & $10$ & $12$ & $14$ & $16$ & $18$ & $20$\\\hline
$b_5$ & $1$ & & & & $1+\omega$ & $1$ & & $-1-\omega$ & & $-1+\omega$\\
$b_6$ & & $1$ & & & $1+\omega$ & $\tfrac23-\tfrac23\omega$ & $-1$ & $-\tfrac23-\tfrac13\omega$ & & $1+2\omega$\\
$b_7$ & & & $1$ & & $-\tfrac13-\tfrac23\omega$ & $-\tfrac23-\tfrac13\omega$ & $1+\omega$ & & $1+\omega$ & $-\omega$\\
$b_8$ & & & & $1$ & $1+\omega$ & $1$ & $-2-\omega$ & $-1-\omega$ & & $3\omega$
\end{tabular}
\end{center}
Each satisfies the cocycle relation $\sum_iQ_iu_i=0$ of Lemma~\ref{lem:cocycle}; for $b_1,\dots,b_4$ this is $v_k+s_kv_k=0$. The same basis, with every value $u_i$ written out, is in the supplementary file \texttt{g34\_data.json}. Then $I=\langle b_1,\dots,b_4\rangle$.

\emph{The matrices.} With $\mathcal S(x,y)=\bar x^{T}S\,y$, the Gram matrix of the skew-Hermitian form is
\[
S=\left(\begin{smallmatrix}0&0&0&0&3&0&0&0\\0&0&0&0&0&1&0&0\\0&0&0&0&0&0&1&0\\0&0&0&0&0&0&0&3\\-3&0&0&0&-3-6\omega&-3-3\omega&4+5\omega&-6-6\omega\\0&-1&0&0&-3\omega&-1-2\omega&1+4\omega&-2-4\omega\\0&0&-1&0&1+5\omega&3+4\omega&-2-4\omega&3+6\omega\\0&0&0&-3&-6\omega&-2-4\omega&3+6\omega&-3-6\omega\end{smallmatrix}\right)
\]
In this basis every $N_j$ has the block form $N_j=\left(\begin{smallmatrix}0&X_j\\0&0\end{smallmatrix}\right)$, with the following $4\times4$ blocks over $K$, numbered as in the table:
{\footnotesize
\[
X_{1}=\left(\begin{smallmatrix}2&1/3-1/3\omega&1/3+2/3\omega&-\omega\\2+\omega&2&0&2+\omega\\-1-2\omega&0&2&-1+\omega\\1+\omega&1/3-1/3\omega&-2/3-1/3\omega&2\end{smallmatrix}\right),\quad X_{2}=\left(\begin{smallmatrix}2&-2\omega&-2/3+2/3\omega&2-2\omega\\6+6\omega&6&-4-2\omega&12+6\omega\\-4-2\omega&-2+2\omega&2&-6\\4+2\omega&2-2\omega&-2&6\end{smallmatrix}\right)
\]
\[
X_{3}=\left(\begin{smallmatrix}0&0&0&0\\0&3&-\omega&3+3\omega\\0&1+\omega&3&3\omega\\0&-\omega&-1-\omega&3\end{smallmatrix}\right),\quad X_{4}=\left(\begin{smallmatrix}0&0&0&0\\0&2&0&2+\omega\\0&0&2&-1+\omega\\0&1/3-1/3\omega&-2/3-1/3\omega&2\end{smallmatrix}\right)
\]
\[
X_{5}=\left(\begin{smallmatrix}0&0&0&0\\0&0&0&0\\0&0&4&-2+2\omega\\0&0&-4/3-2/3\omega&3\end{smallmatrix}\right),\quad X_{6}=\left(\begin{smallmatrix}0&0&0&0\\0&0&0&0\\0&0&2&0\\0&0&0&0\end{smallmatrix}\right)
\]
\[
X_{7}=\left(\begin{smallmatrix}0&0&0&0\\0&0&0&0\\0&0&2&-1+\omega\\0&0&-2/3-1/3\omega&2\end{smallmatrix}\right),\quad X_{8}=\left(\begin{smallmatrix}0&0&0&0\\0&0&0&0\\0&0&0&0\\0&0&0&2\end{smallmatrix}\right)
\]
\[
X_{9}=\left(\begin{smallmatrix}2&2&-4/3-2/3\omega&2\\6&6&-4-2\omega&6\\-2+2\omega&-2+2\omega&2&-2+2\omega\\2&2&-4/3-2/3\omega&2\end{smallmatrix}\right),\quad X_{10}=\left(\begin{smallmatrix}3&2-\omega&-2-\omega&3\\9+3\omega&8&-5-4\omega&9+3\omega\\-3+3\omega&-1+4\omega&3&-3+3\omega\\3&2-\omega&-2-\omega&3\end{smallmatrix}\right)
\]
\[
X_{11}=\left(\begin{smallmatrix}5&4/3-7/3\omega&-4/3+1/3\omega&3-2\omega\\11+7\omega&8&-5-3\omega&11+4\omega\\-5-\omega&-2+3\omega&4&-5+2\omega\\5+2\omega&7/3-4/3\omega&-7/3-2/3\omega&5\end{smallmatrix}\right),\quad X_{12}=\left(\begin{smallmatrix}7&2/3-11/3\omega&-2/3+5/3\omega&3-4\omega\\13+11\omega&9&-5-2\omega&13+5\omega\\-7-5\omega&-3+2\omega&5&-7+\omega\\7+4\omega&8/3-5/3\omega&-8/3-1/3\omega&7\end{smallmatrix}\right)
\]
\[
X_{13}=\left(\begin{smallmatrix}5&4/3-7/3\omega&-1/3+4/3\omega&3-2\omega\\11+7\omega&8&-3&11+4\omega\\-5-4\omega&-3&3&-5-\omega\\5+2\omega&7/3-4/3\omega&-4/3+1/3\omega&5\end{smallmatrix}\right),\quad X_{14}=\left(\begin{smallmatrix}2&4/3-4/3\omega&-4/3-2/3\omega&2\\8+4\omega&8&-4-4\omega&8+4\omega\\-2+2\omega&4\omega&2&-2+2\omega\\2&4/3-4/3\omega&-4/3-2/3\omega&2\end{smallmatrix}\right)
\]
\[
X_{15}=\left(\begin{smallmatrix}0&0&0&0\\0&2&-2-2\omega&4+2\omega\\0&2\omega&2&-2+2\omega\\0&2/3-2/3\omega&-4/3-2/3\omega&2\end{smallmatrix}\right),\quad X_{16}=\left(\begin{smallmatrix}2&-\omega&1/3+2/3\omega&1-\omega\\3+3\omega&3&-2-\omega&6+3\omega\\-1-2\omega&-1+\omega&2&-3\\2+\omega&1-\omega&-1&3\end{smallmatrix}\right)
\]
\[
X_{17}=\left(\begin{smallmatrix}2&2/3-2/3\omega&0&2\\4+2\omega&2&0&4+2\omega\\0&0&0&0\\2&2/3-2/3\omega&0&2\end{smallmatrix}\right)
\]
}
So (H1) is the block form itself, and (H2) holds because the upper left $4\times4$ block of $S$ vanishes. On $\Vp/I=\langle b_5,\dots,b_8\rangle$ the form $\beta_j$ has matrix $DX_j$, where $D=\operatorname{diag}(-3,-1,-1,-3)$ is the lower left block of $S$. So (H3) and (H4) are statements about $4\times4$ matrices that can be checked by hand: each $DX_j$ is Hermitian and negative semidefinite, of the rank in the table; $\sum_jDX_j$ is negative definite; and the $X_j$ span a $\Q$-space of dimension $16$ in the $32$-dimensional $\Q$-space of $4\times4$ matrices over $K$. For example $DX_1$ has first row $(-6,\,-1+\omega,\,-1-2\omega,\,3\omega)$. The certificate has two levels. \emph{The printed level:} from $S$ and the $X_j$ alone, (H1)--(H4) are statements about $4\times4$ matrices over $K$; the supplementary script \texttt{check\_g34.py} (Python~3, exact arithmetic over $K$, no other dependencies) checks them, and \texttt{python3 check\_g34.py} ends with the line \texttt{ALL (H1)-(H4) HOLD: True} after printing the ranks and signatures of the table. \emph{The derivation level:} from the printed roots, tuple, curves and basis $b_1,\dots,b_8$, the matrices $S$ and $X_j$ are recomputed by the procedure described below; the supplementary files contain two separately written programs that do this, with their outputs and checksums.

\emph{How to recompute.} $\Vp=H^1_{\rm par}(U_C,L')$ is the space of parabolic cocycles modulo coboundaries for the free group on the twenty loops with monodromy $\sigma(\mathbf g)$, of $K$-dimension $8$; $\mathcal S$ is given by Lemma~\ref{lem:cocycle}; each $T_j$ is the $k_j$-th iterate of the pullback formula of Section~\ref{sec:boundary}, and $N_j=\log T_j$ since every $m_j=1$. All of this is exact linear algebra over $K$ in dimension $8$. The supplementary files contain two separately written programs that compute the table from the tuple and agree, and the explicit supplier identity of Remark~\ref{rem:witness} with its rebuild script.

%% ------------------------------------------------------------------
\begin{thebibliography}{99}
\bibitem{Abdulali} S. Abdulali, \emph{Abelian varieties of type III and the Hodge conjecture}, Internat. J. Math. 10 (1999), 667--675.
\bibitem{ADFIO} V. Alexeev, R. Donagi, G. Farkas, E. Izadi and A. Ortega, \emph{The uniformization of the moduli space of principally polarized abelian 6-folds}, J. Reine Angew. Math. 761 (2020), 163--217.
\bibitem{BL} C. Birkenhake and H. Lange, \emph{Complex Abelian Varieties}, second edition, Grundlehren 302, Springer, 2004.
\bibitem{Brosnan} P. Brosnan, \emph{Discriminants of Hermitian forms, maximal degenerations and Kontsevich's tropical approach to the Hodge conjecture}, arXiv:2609.14169, 2026.
\bibitem{CKS} E. Cattani, A. Kaplan and W. Schmid, \emph{Degeneration of Hodge structures}, Ann. of Math. 123 (1986), 457--535.
\bibitem{Clemens} C. H. Clemens, \emph{Degeneration of K\"ahler manifolds}, Duke Math. J. 44 (1977), 215--290.
\bibitem{DW1} M. Dettweiler and S. Wewers, \emph{Variation of local systems and parabolic cohomology}, Israel J. Math. 156 (2006), 157--185.
\bibitem{DW2} M. Dettweiler and S. Wewers, \emph{Variation of parabolic cohomology and Poincar\'e duality}, arXiv:math/0411119.
\bibitem{FGP} P. Frediani, A. Ghigi and M. Penegini, \emph{Shimura varieties in the Torelli locus via Galois coverings}, Int. Math. Res. Not. IMRN 2015, 10595--10623.
\bibitem{FriedVolklein} M. Fried and H. V\"olklein, \emph{The inverse Galois problem and rational points on moduli spaces}, Math. Ann. 290 (1991), 771--800.
\bibitem{HarrisMumford} J. Harris and D. Mumford, \emph{On the Kodaira dimension of the moduli space of curves}, Invent. Math. 67 (1982), 23--86.
\bibitem{Hatcher} A. Hatcher, \emph{Algebraic Topology}, Cambridge University Press, 2002.
\bibitem{Howe} R. Howe, \emph{Remarks on classical invariant theory}, Trans. Amer. Math. Soc. 313 (1989), 539--570.
\bibitem{Knudsen} F. Knudsen, \emph{The projectivity of the moduli space of stable curves II: the stacks $M_{g,n}$}, Math. Scand. 52 (1983), 161--199.
\bibitem{Koike} K. Koike, \emph{Algebraicity of some Weil Hodge classes}, Canad. Math. Bull. 47 (2004), 566--572.
\bibitem{KraftProcesi} H. Kraft and C. Procesi, \emph{Classical Invariant Theory: A Primer}, preliminary version, July 1996, University of Basel.
\bibitem{Markman19} E. Markman, \emph{The monodromy of generalized Kummer varieties and algebraic cycles on their intermediate Jacobians}, arXiv:1805.11574.
\bibitem{Markman25} E. Markman, \emph{Cycles on abelian $2n$-folds of Weil type from secant sheaves on abelian $n$-folds}, arXiv:2502.03415.
\bibitem{Markman25b} E. Markman, \emph{Secant sheaves and Weil classes on abelian varieties}, arXiv:2509.23403.
\bibitem{Milne} J. S. Milne, \emph{Lefschetz classes on abelian varieties}, Duke Math. J. 96 (1999), 639--675.
\bibitem{Mostaed} A. Mostaed, \emph{McMullen's curve, the Weil locus, and the Hodge conjecture for abelian sixfolds}, arXiv:2603.20268, 2026.
\bibitem{Schmid} W. Schmid, \emph{Variation of Hodge structure: the singularities of the period mapping}, Invent. Math. 22 (1973), 211--319.
\bibitem{Schoen88} C. Schoen, \emph{Hodge classes on self-products of a variety with an automorphism}, Compositio Math. 65 (1988), 3--32.
\bibitem{Schoen98} C. Schoen, \emph{Addendum to: Hodge classes on self-products of a variety with an automorphism}, Compositio Math. 114 (1998), 329--336.
\bibitem{Schapira} P. Schapira, \emph{An Introduction to Categories and Sheaves}, lecture notes, version of 1 October 2023.
\bibitem{SerreLin} J.-P. Serre, \emph{Linear Representations of Finite Groups}, Graduate Texts in Mathematics 42, Springer, 1977.
\bibitem{SevernG32} B. Severn, \emph{Algebraicity of Weil classes on abelian sixfolds of Weil type for $\Q(\sqrt{-3})$ with discriminant $-2$}, preprint, version 0.10, HexagonMath hexagon:2610.00246v1, 2026.
\bibitem{ShephardTodd} G. C. Shephard and J. A. Todd, \emph{Finite unitary reflection groups}, Canad. J. Math. 6 (1954), 274--304.
\bibitem{vanGeemen} B. van Geemen, \emph{An introduction to the Hodge conjecture for abelian varieties}, in: Algebraic Cycles and Hodge Theory (Torino, 1993), Lecture Notes in Math. 1594, Springer, 1994, 233--252.
\bibitem{vGV} B. van Geemen and A. Verra, \emph{Quaternionic Pryms and Hodge classes}, Topology 42 (2003), 35--53.
\bibitem{Voight} J. Voight, \emph{Quaternion Algebras}, Graduate Texts in Mathematics 288, Springer, 2021.
\bibitem{Voisin} C. Voisin, \emph{La conjecture de Hodge pour les vari\'et\'es ab\'eliennes de dimension au plus 5 [d'apr\`es Markman]}, S\'eminaire Bourbaki, expos\'e 1248, 2026.
\bibitem{VoisinII} C. Voisin, \emph{Hodge Theory and Complex Algebraic Geometry II}, Cambridge Studies in Advanced Mathematics 77, Cambridge University Press, 2003.
\bibitem{Weil} A. Weil, \emph{Abelian varieties and the Hodge ring}, in: Collected Papers III, Springer, 1980, 421--429.
\bibitem{Weyl} H. Weyl, \emph{The Classical Groups: Their Invariants and Representations}, second edition, Princeton University Press, 1946.
\end{thebibliography}
\end{document}
